% Source: 20.determinants.a/.b/.c, which number from 20.a.1, 20.b.1 and 20.c.1
% respectively. The prefix is re-declared at each \section below.
\renewcommand{\thetheorem}{20.a.\arabic{theorem}}
\setcounter{theorem}{0}
\setcounter{chapter}{19} % Chapter 20
% ---------------------------------------------------------------------------
% Provenance of the prose in this file.
%   % Extractor: <model>   the block below was transcribed from Prof. Biran's
%                          handwritten notes by that model
%   % Creator:   <model>   the block below was written by that model and has no
%                          counterpart in the notes
% A tag holds until the next tag.
% ---------------------------------------------------------------------------

\chapter{Determinants}

% Extractor: Gemini 3.1 Pro
\section{Introduction and Motivation}

\subsection{The special case of \texorpdfstring{$2 \times 2$}{2x2} matrices}
In linear algebra, determining whether a matrix is invertible is a fundamental problem. Let us begin by analyzing the familiar $2 \times 2$ case. Let $A = (\begin{smallmatrix} a & b \\ c & d \end{smallmatrix}) \in \M_{2 \times 2}(K)$. There is a very simple, well-known algebraic criterion to check whether $A$ is invertible.

\begin{proposition}[Invertibility of $2 \times 2$ Matrices]\label{prop:inv_2x2}
  The matrix $A$ is invertible if and only if $ad - bc \neq 0$. Moreover, if $A$ is invertible, then its inverse is given by the explicit formula:
  \[
  A^{-1} = \frac{1}{ad - bc} \begin{pmatrix} d & -b \\ -c & a \end{pmatrix}.
  \]
  The scalar quantity $ad - bc$ is called the \newterm{determinant} of $A$, denoted $\det(A)$.
\end{proposition}

\begin{proof}
  First, $A$ is not invertible if and only if its rank is strictly less than $2$; this is \cref{aiprop:invertible_iff_full_rank}, read in the contrapositive. Because $A$ is a $2 \times 2$ matrix, having a rank less than 2 implies that its column vectors are linearly dependent: the column rank equals the rank by \cref{thm:rank_equality}, so two columns spanning a space of dimension at most $1$ cannot be independent, by \cref{thm:wrong_length_lists} \textbf{(2)}.

% Extractor: Opus 5
  \begin{exercise}[Dependence of Two Columns]\label{exc:lin_dep_2x2}
    Check that the columns of $A$ being linearly dependent means that either the first column is a scalar multiple of the second, $\binom{a}{c} = \lambda \binom{b}{d}$ for some $\lambda \in K$, or the second is a scalar multiple of the first, $\binom{b}{d} = \tau \binom{a}{c}$ for some $\tau \in K$.
  \exinfo{This exercise is taken from Prof.~Biran's handwritten lecture notes.}
  \exsol{sol:lin_dep_2x2}
\end{exercise}

% Extractor: Gemini 3.1 Pro

  If $\binom{a}{c} = \lambda \binom{b}{d}$, then $a = \lambda b$ and $c = \lambda d$. Consequently, $ad - bc = (\lambda b)d - b(\lambda d) = 0$.
  If, on the other hand, $\binom{b}{d} = \tau \binom{a}{c}$, then $b = \tau a$ and $d = \tau c$. Consequently, $ad - bc = a(\tau c) - (\tau a)c = 0$.
  Thus, we have proved that if $A$ is not invertible, then $ad - bc = 0$.

  Conversely, assume $ad - bc = 0$. If $a = 0$, then $bc = 0$, so either $b = 0$ or $c = 0$. This means either $A = \left(\begin{smallmatrix} 0 & 0 \\ c & d \end{smallmatrix}\right)$ or $A = \left(\begin{smallmatrix} 0 & b \\ 0 & d \end{smallmatrix}\right)$. In both of these cases, $A$, clearly, has linearly dependent columns and is therefore not invertible.

  Now, assume $a \neq 0$. This implies $d = \frac{bc}{a}$, and thus we can write $A = \left(\begin{smallmatrix} a & b \\ c & \frac{bc}{a} \end{smallmatrix}\right)$. We can observe that the second column of $A$ is exactly a multiple of the first column because $\binom{b}{\frac{bc}{a}} = \frac{b}{a} \binom{a}{c}$. This implies that the columns of $A$ are linearly dependent. Hence, $\rank(A) < 2$, which implies $A$ is not invertible.
\end{proof}

To confirm the sufficiency of this condition and to find the explicit inverse, the formula for $A^{-1}$ can be verified by direct matrix multiplication:
\[
\frac{1}{ad-bc} \begin{pmatrix} d & -b \\ -c & a \end{pmatrix} \begin{pmatrix} a & b \\ c & d \end{pmatrix} = \frac{1}{ad-bc} \begin{pmatrix} ad-bc & bd-bd \\ -ac+ac & -bc+ad \end{pmatrix} = \begin{pmatrix} 1 & 0 \\ 0 & 1 \end{pmatrix}.
\]
An inverse has to satisfy the identity in both orders, and the product taken the other way round is the same computation:
\[
\frac{1}{ad-bc} \begin{pmatrix} a & b \\ c & d \end{pmatrix} \begin{pmatrix} d & -b \\ -c & a \end{pmatrix} = \frac{1}{ad-bc} \begin{pmatrix} ad-bc & -ab+ab \\ cd-dc & -cb+ad \end{pmatrix} = \begin{pmatrix} 1 & 0 \\ 0 & 1 \end{pmatrix}.
\]

\begin{exercise}[Invertibility and Nilpotency of $2 \times 2$ Matrices \faHeart]
\label{exc:invertibility_nilpotency_2x2}
Let $A = \begin{pmatrix} a & b \\ c & d \end{pmatrix} \in \M_{2 \times 2}(K)$ be a $2 \times 2$ matrix over a field $K$.
\begin{enumerate}
    \item Show that $A$ is invertible if and only if $ad - bc \neq 0$.
    \item Show that
    \[
    A^2 = 0 \quad \Longleftrightarrow \quad \det(tI_2 - A) = t^2
    \]
    as polynomials in the polynomial ring $K[t]$.
\end{enumerate}
\exhint[Hint (Y. St\"aubli)]{(a) For the forward implication, use the explicit adjugate formula for $A^{-1}$; for the reverse implication, use the multiplicativity of the determinant $\det(AB) = \det(A)\det(B)$. (b) First show that $\det(tI_2 - A) = t^2 - (\operatorname{tr} A)t + \det A$, where $\operatorname{tr}(A) = a + d$, and compare coefficients.}
\exinfo{This exercise is Problem 2 of Problem Sheet 16, Linear Algebra II, Spring Semester 2026.}
\exsol{sol:invertibility_nilpotency_2x2}
\end{exercise}

\begin{goals}\label{goals:determinants}
  While the $2 \times 2$ formula is easy to memorize, it does not immediately suggest how to handle larger matrices. Our main goals now are to make the determinant conceptually rigorous and to generalize the determinant function to arbitrary $n \times n$ matrices.
\end{goals}

\subsection{Multilinear and Alternating Functions}

We wish to generalize the determinant to arbitrary dimensions. To this end, we should view an $n \times n$ matrix not just as a grid of numbers, but as a collection of $n$ row vectors.

\begin{definition}[$n$-Linearity]\label{def:n_linear}
  Let $D : \M_{n \times n}(K) \to K$ be a function. We say that $D$ is \newterm{$n$-linear} if for every $1 \le i \le n$, $D$ is linear as a function of the $i$-th row when all the other $n-1$ rows are held fixed.
\end{definition}

We can think of the matrix space $\M_{n \times n}(K)$ as a Cartesian product of row vectors: $K^n_{\text{row}} \times \dots \times K^n_{\text{row}}$. If $A$ has rows $\alpha_1, \dots, \alpha_n$, we write $D(A) = D(\alpha_1, \dots, \alpha_n)$.
Saying that $D$ is $n$-linear means that for any chosen row index $i$, the following two linearity conditions hold:
\begin{enumerate}
  \item $D(\alpha_1, \dots, \alpha_i + \alpha_i', \dots, \alpha_n) = D(\alpha_1, \dots, \alpha_i, \dots, \alpha_n) + D(\alpha_1, \dots, \alpha_i', \dots, \alpha_n)$.
  \item $D(\alpha_1, \dots, c \cdot \alpha_i, \dots, \alpha_n) = c \cdot D(\alpha_1, \dots, \alpha_i, \dots, \alpha_n)$ for all scalars $c \in K$.
\end{enumerate}

\begin{example}[Product of Row Entries]\label{ex:n_linear_product}
  For a matrix $A \in \M_{n \times n}(K)$, write $A(i,j)$ for its entry at place $i,j$. Let $k_1, \dots, k_n$ be a fixed sequence of integers such that $1 \le k_i \le n$ for all $i$. Fix also a scalar $c \in K$. Consider $D : \M_{n \times n}(K) \to K$ defined by the product of one entry from each row:
  \[
  D(A) := c \cdot A(1, k_1) \cdot A(2, k_2) \cdots A(n, k_n).
  \]
  We claim that $D$ is an $n$-linear function.

  To see this, consider $D(\alpha_1, \dots, \alpha_i, \dots, \alpha_n) = b \cdot A(i, k_i)$, where $b \in K$ is the product of all the fixed terms. Notice that $b$ depends only on the rows $\alpha_1, \dots, \alpha_{i-1}, \alpha_{i+1}, \dots, \alpha_n$ and does not depend on row $i$.
  Let $\alpha'_i = (\alpha'_{i,1}, \dots, \alpha'_{i,n})$ be another possible $i$-th row. Then, by standard distribution:
  \[
  D(\dots, \alpha_i + \alpha'_i, \dots) = b \cdot (A(i, k_i) + \alpha'_{i, k_i}) = D(\dots, \alpha_i, \dots) + D(\dots, \alpha'_i, \dots).
  \]
  Similarly, for any scalar $\lambda \in K$:
  \[
  D(\dots, \lambda \cdot \alpha_i, \dots) = b \cdot (\lambda A(i, k_i)) = \lambda \cdot D(\dots, \alpha_i, \dots).
  \]
  Therefore, $D$ is indeed $n$-linear.
\end{example}

\begin{exercise}[Monomial Matrix Functions and Multilinearity \faHeart]
\label{exc:monomial_matrix_multilinear}
Let $K$ be a subfield of the complex numbers $\mathbb{C}$, and let $n$ be a positive integer. Let $j_1, \dots, j_n$ and $k_1, \dots, k_n$ be positive integers in $\{1, \dots, n\}$. For each $A \in \M_{n \times n}(K)$, define
\[
D(A) := A(j_1, k_1) A(j_2, k_2) \cdots A(j_n, k_n).
\]
Prove that $D$ is $n$-linear (as a function of the rows of $A$) if and only if the row indices $j_1, \dots, j_n$ are pairwise distinct.
\exhint[Hint (Y. St\"aubli)]{If two row indices are equal, say $j_1 = j_2$, show that scaling that row by $c \in K$ scales the product by $c^2$ rather than $c$ (explain why $K \subseteq \mathbb{C}$ prevents $c^2 = c$ for all $c$, in contrast to finite fields like $\mathbb{F}_2$). If all $j_i$ are distinct, each row appears exactly once in the product, so $n$-linearity holds directly.}
\exinfo{This exercise is Problem 4 of Problem Sheet 16, Linear Algebra II, Spring Semester 2026.}
\exsol{sol:monomial_matrix_multilinear}
\end{exercise}

\begin{example}[Diagonal Product]\label{ex:diagonal_product}
  A more concrete example: $D(A) := A_{11} \cdot A_{22} \cdot \dots \cdot A_{nn}$ (i.e., the product of the terms on the main diagonal) is an $n$-linear function.
\end{example}

Now, let us try to find all 2-linear functions $D: \M_{2 \times 2}(K) \to K$. Write the identity matrix as $I = \left(\begin{smallmatrix} -\epsilon_1- \\ -\epsilon_2- \end{smallmatrix}\right)$, where $\epsilon_1 = (1, 0)$ and $\epsilon_2 = (0, 1)$.

If $D$ is 2-linear, then we have:
\begin{align*}
  D(A) &= D(A_{11}\epsilon_1 + A_{12}\epsilon_2, A_{21}\epsilon_1 + A_{22}\epsilon_2) \\
  &= A_{11} D(\epsilon_1, A_{21}\epsilon_1 + A_{22}\epsilon_2) + A_{12} D(\epsilon_2, A_{21}\epsilon_1 + A_{22}\epsilon_2) \\
  &= A_{11}A_{21}D(\epsilon_1, \epsilon_1) + A_{11}A_{22}D(\epsilon_1, \epsilon_2) + A_{12}A_{21}D(\epsilon_2, \epsilon_1) + A_{12}A_{22}D(\epsilon_2, \epsilon_2).
\end{align*}
Therefore, $D$ is completely determined by its values on the following four pairings: $D(\epsilon_1, \epsilon_1)$, $D(\epsilon_1, \epsilon_2)$, $D(\epsilon_2, \epsilon_1)$, and $D(\epsilon_2, \epsilon_2)$.

\begin{lemma}[Closure of $n$-Linear Functions]\label{lem:linear_comb_n_linear}
  Any linear combination of $n$-linear functions is again an $n$-linear function.
\end{lemma}

\begin{proof}
  It is enough to prove this for a linear combination of two $n$-linear functions. Let $D, E : \M_{n \times n}(K) \to K$ be $n$-linear functions and $a, b \in K$. We define the combined function as $(aD + bE)(A) := aD(A) + bE(A)$.
  By expanding $(aD+bE)(\dots, \alpha_i + \alpha'_i, \dots)$ and $(aD+bE)(\dots, \lambda \alpha_i, \dots)$ using the individual linearities of $D$ and $E$, we see the result naturally holds.
\end{proof}

\begin{exercise}[Multilinear Candidate Functions on $3 \times 3$ Matrices \faHeart]
\label{exc:multilinear_candidates_3x3}
Each of the following expressions defines a function $D \colon \M_{3 \times 3}(\mathbb{R}) \to \mathbb{R}$. In which of these cases is $D$ a $3$-linear function of the rows of the matrix?
\begin{enumerate}
    \item $D(A) = A_{11} + A_{22} + A_{33}$;
    \item $D(A) = A_{11}^2 + 3A_{11} A_{22}$;
    \item $D(A) = A_{11} A_{12} A_{33}$;
    \item $D(A) = A_{13} A_{22} A_{32} + 5A_{12} A_{22} A_{32}$;
    \item $D(A) = 0$;
    \item $D(A) = 1$.
\end{enumerate}
\exhint[Hint (Y. St\"aubli)]{Here, $3$-linear with respect to the rows is meant. For those functions that fail to be $3$-linear, provide an explicit counterexample to linearity in some row.}
\exinfo{This exercise is Problem 3 of Problem Sheet 16, Linear Algebra II, Spring Semester 2026.}
\exsol{sol:multilinear_candidates_3x3}
\end{exercise}

\begin{example}[Two-by-Two Determinant Function]\label{ex:2_linear_det}
  Consider $D : \M_{2 \times 2}(K) \to K$ defined by $D(A) := A_{11}A_{22} - A_{12}A_{21}$. The function $D$ is $2$-linear because it can be written as $D = D_1 + D_2$, where $D_1(A) = A_{11}A_{22}$ and $D_2(A) = -A_{12}A_{21}$, both of which are 2-linear by our previous example. Furthermore, $D(I) = 1$, and if $A'$ is obtained from $A$ by interchanging the rows, then $D(A') = -D(A)$.

  Indeed, $D(A') = D\left(\begin{smallmatrix} A_{21} & A_{22} \\ A_{11} & A_{12} \end{smallmatrix}\right) = A_{21} \cdot A_{12} - A_{22}A_{11} = -(A_{11}A_{22} - A_{12}A_{21}) = -D(A)$.
\end{example}

\begin{exercise}[The $2 \times 2$ Determinant as an Alternating $2$-Linear Form \faHeart]
\label{exc:det_2x2_alternating_bilinear}
Show that the determinant of a $2 \times 2$ matrix, $\det \begin{pmatrix} a & b \\ c & d \end{pmatrix} = ad - bc$, is an alternating $2$-linear function with respect to both its rows and its columns.
\exhint[Hint (Y. St\"aubli)]{Using the explicit formula $\det A = ad - bc$, verify that $\det$ is linear in each column (holding the other fixed) and linear in each row (holding the other fixed), and show that $\det A = 0$ whenever the two columns (or the two rows) are equal.}
\exinfo{This exercise is Problem 1 of Problem Sheet 16, Linear Algebra II, Spring Semester 2026.}
\exsol{sol:det_2x2_alternating_bilinear}
\end{exercise}

\begin{definition}[Alternating Function]\label{def:alternating_func}
  Let $D : \M_{n \times n}(K) \to K$ be an $n$-linear function. We say that $D$ is \newterm{alternating} if for every matrix $A$ in which some two rows are equal, we have $D(A) = 0$.
\end{definition}

\begin{proposition}[Adjacent Swap Property]\label{prop:alternating_swap}
  Let $D : \M_{n \times n}(K) \to K$ be an $n$-linear function, and assume that $D$ has the property that whenever $A$ has two equal \emph{adjacent} rows, then $D(A) = 0$. Then both of the following hold:
  \begin{enumerate}[label=\textbf{(\alph*)}]
    \item $D$ is alternating, that is, $D(A) = 0$ whenever \emph{any} two rows of $A$ are equal, adjacent or not.
    \item For any matrix $B$, if $B'$ is obtained from $B$ by interchanging two of its rows, then $D(B') = -D(B)$.
  \end{enumerate}
\end{proposition}

\begin{proof}
  We begin with the proof of \textbf{(b)}. Assume first that $B'$ is obtained from $B$ by interchanging two adjacent rows (say, row $k$ and row $k+1$).
  Consider evaluating $D$ on a matrix where both the $k$-th and $(k+1)$-th rows are equal to $\beta_k + \beta_{k+1}$. By our assumption, this yields $0$.
  Expanding this via $n$-linearity gives:
  \begin{align*}
    D(\dots, \beta_k + \beta_{k+1}, \beta_k + \beta_{k+1}, \dots)
      &= \underbrace{D(\dots, \beta_k, \beta_k, \dots)}_{=\, 0}
       + D(\dots, \beta_k, \beta_{k+1}, \dots) \\
      &\quad + D(\dots, \beta_{k+1}, \beta_k, \dots)
       + \underbrace{D(\dots, \beta_{k+1}, \beta_{k+1}, \dots)}_{=\, 0} \;=\; 0.
  \end{align*}
  The two braced terms carry two equal adjacent rows, so they vanish by assumption, and the identity collapses to
  \[
  \underbrace{D(\dots, \beta_k, \beta_{k+1}, \dots)}_{D(B)}
  + \underbrace{D(\dots, \beta_{k+1}, \beta_k, \dots)}_{D(B')} = 0,
  \]
  which naturally implies $D(B') = -D(B)$.

  Now suppose $B'$ is obtained from $B$ by interchanging rows $k$ and $\ell$ where $k < \ell$. We can obtain $B'$ by doing a sequence of adjacent row interchanges. Moving row $\ell$ to position $k$ requires $r = \ell-k$ adjacent interchanges. Next, moving the original row $k$ down to position $\ell$ requires $r-1$ adjacent interchanges. The total number of interchanges is $2r - 1$. Since this is always an odd number, we flip the sign an odd number of times. Thus, $D(B') = (-1)^{2r-1} D(B) = -D(B)$.

  We now prove \textbf{(a)}. Let $A$ be a matrix in which row $i$ equals row $j$ where $i < j$. If $j = i+1$, then by assumption $D(A) = 0$. If $j > i+1$, we interchange row $j$ with row $i+1$ to obtain a matrix $A'$ with two equal adjacent rows.
  By assumption $D(A') = 0$, but by \textbf{(b)}, $D(A') = -D(A)$ (since we performed an odd number of adjacent swaps), which implies $D(A) = 0$.
\end{proof}

\subsection{The Determinant Function}

\begin{definition}[Determinant Function]\label{def:determinant_func}
  A function $D : \M_{n \times n}(K) \to K$ is called a \qt{determinant function} if it is $n$-linear, it is alternating, and it normalizes the identity matrix such that $D(I) = 1$.
\end{definition}

\begin{example}[Warmup]
  Let us find all determinant functions on $M_{2 \times 2}(K)$. Write $I = \left(\begin{smallmatrix} 1 & 0 \\ 0 & 1 \end{smallmatrix}\right) = \left(\begin{smallmatrix} -\epsilon_1- \\ -\epsilon_2- \end{smallmatrix}\right)$, meaning $\epsilon_1 = (1,0)$ and $\epsilon_2 = (0,1)$.
  If $D$ is a $2$-linear function, we can expand it algebraically:
  \begin{align*}
    D(A) &= D(A_{11}\epsilon_1 + A_{12}\epsilon_2, A_{21}\epsilon_1 + A_{22}\epsilon_2) \\
    &= A_{11} D(\epsilon_1, A_{21}\epsilon_1 + A_{22}\epsilon_2) + A_{12} D(\epsilon_2, A_{21}\epsilon_1 + A_{22}\epsilon_2) \\
    &= A_{11}A_{21}D(\epsilon_1, \epsilon_1) + A_{11}A_{22}D(\epsilon_1, \epsilon_2) \\
    &\quad + A_{12}A_{21}D(\epsilon_2, \epsilon_1) + A_{12}A_{22}D(\epsilon_2, \epsilon_2).
  \end{align*}
  Consequently, $D$ is completely determined by its values on these standard basis row pairings.

  If $D$ is alternating, then evaluating equal rows gives zero, meaning $D(\epsilon_1, \epsilon_1) = D(\epsilon_2, \epsilon_2) = 0$. Additionally, swapping rows flips the sign, so $D(\epsilon_2, \epsilon_1) = -D(\epsilon_1, \epsilon_2) = -D(I)$.
  Finally, if $D$ is a determinant function, $D(I) = 1$. Substituting these values back into our expansion, we perfectly recover the familiar formula:
  \[
  D(A) = A_{11}A_{22}(1) + A_{12}A_{21}(-1) = A_{11}A_{22} - A_{12}A_{21}.
  \]
  Therefore, the $2 \times 2$ determinant function is completely unique!
\end{example}

\begin{definition}[Submatrix]\label{def:minor_submatrix}
  Let $A \in \M_{n \times n}(K)$, $n \ge 2$, and $1 \le i \le n$, $1 \le j \le n$. We denote by $A(i|j) \in \M_{(n-1) \times (n-1)}(K)$ the \qt{submatrix} obtained from $A$ by deleting row $i$ and column $j$. If $D$ is an $(n-1)$-linear function, we intuitively define $D_{ij}(A) := D(A(i|j))$.
\end{definition}

\begin{theorem}[Recursive Construction of Determinants]\label{thm:recursive_det}
  Let $n \ge 2$, and let $D$ be an $(n-1)$-linear function which is alternating. Let $1 \le j \le n$. Define $E_j : \M_{n \times n}(K) \to K$ by
  \[
  E_j(A) := \sum_{i=1}^n (-1)^{i+j} A_{ij} D_{ij}(A).
  \]
  Then, $E_j$ is an $n$-linear function and it is alternating. Furthermore, if $D$ is a determinant function, then so is $E_j$.
\end{theorem}

\begin{proof}
  Clearly, $D_{ij}(A)$ is independent of the $i$-th row of $A$ because that row is deleted in the construction of $A(i|j)$. Since $D$ is $(n-1)$-linear, the map $A \mapsto D_{ij}(A)$ is linear when viewed as a function of each of its remaining rows. The function $A \mapsto A_{ij} D_{ij}(A)$ is therefore $n$-linear. Since linear combinations of $n$-linear functions are $n$-linear, $E_j(A)$ is $n$-linear too.

  We show that $E_j$ is alternating. By \Cref{prop:alternating_swap}, it is enough to show $E_j(A) = 0$ whenever $A$ has two equal adjacent rows.
  Assume $\alpha_k = \alpha_{k+1}$ for some index $k$. For any $i \neq k, k+1$, the submatrix $A(i|j)$ still has two equal rows, hence $D_{ij}(A) = 0$. This leaves only two non-zero terms in the sum:
  \[
  E_j(A) = (-1)^{k+j} A_{kj} D_{kj}(A) + (-1)^{k+1+j} A_{(k+1)j} D_{(k+1)j}(A).
  \] % (since all other terms $D_{ij}(A)$ are zero because $A(i|j)$ would still have two identical rows)
  But, since $\alpha_k = \alpha_{k+1}$, we have $A_{kj} = A_{(k+1)j}$ and $A(k|j) = A(k+1|j)$, meaning $D_{kj}(A) = D_{(k+1)j}(A)$. Thus, the two terms perfectly cancel out, yielding $E_j(A) = 0$. This proves $E_j$ is alternating.

  Finally, assume $D$ is a determinant function, meaning $D(I_{n-1}) = 1$. Note that $I_n(j|j) = I_{n-1}$ for all $j$. When computing $E_j(I_n)$, the term $(I_n)_{ij}$ is $1$ if $i=j$, and $0$ if $i \neq j$. Therefore, the sum collapses to a single term:
  \[
  E_j(I_n) = (-1)^{j+j} D_{jj}(I_n) = D(I_n(j|j)) = 1.
  \]
  This shows that $E_j$ is a valid determinant function.
\end{proof}

\begin{corollary}[Existence of Determinants]\label{cor:det_exists}
  For all $n \ge 1$, at least one determinant function exists.
\end{corollary}

% Creator: Opus 5
% Stated without proof in both the notes and the transcript, although the proof
% of \cref{lem:permutation_parity} below rests on it. The induction is three
% lines and its base case is worth stating, since alternating is vacuous for
% n = 1.
\begin{proof}
  By induction on $n$. For $n = 1$, the function $D(A) := A_{11}$ is $1$-linear, and it satisfies $D(I_1) = 1$; it is alternating for the empty reason that a $1 \times 1$ matrix has only one row and so cannot have two equal ones. For the step, suppose a determinant function $D$ on $\M_{(n-1) \times (n-1)}(K)$ exists, with $n \ge 2$. Then \cref{thm:recursive_det}, applied with any fixed $j$, produces the determinant function $E_j$ on $\M_{n \times n}(K)$.
\end{proof}

\begin{example}[Three-by-Three Determinant Expansion]
  Consider $A \in \M_{3 \times 3}(K)$. Using the uniquely defined $2 \times 2$ determinant function $|B| := B_{11}B_{22} - B_{12}B_{21}$, we can explicitly write out $E_1(A)$ (expanding along the first column) as:
  \[
  E_1(A) = A_{11} \begin{vmatrix} A_{22} & A_{23} \\ A_{32} & A_{33} \end{vmatrix} - A_{21} \begin{vmatrix} A_{12} & A_{13} \\ A_{32} & A_{33} \end{vmatrix} + A_{31} \begin{vmatrix} A_{12} & A_{13} \\ A_{22} & A_{23} \end{vmatrix}.
  \]
  We can similarly construct $E_2(A)$ and $E_3(A)$ along the other columns, and all of these form valid determinant functions on $M_{3 \times 3}(K)$.
\end{example}

\begin{exercise}[Laplace Expansion for $3 \times 3$ Matrices]
\label{exc:laplace_expansion_3x3_alternating_trilinear}
Let $K$ be a field. Consider the function $D \colon \M_{3 \times 3}(K) \to K$ defined by
\[
D(A) := A_{11} \det \begin{pmatrix} A_{22} & A_{23} \\ A_{32} & A_{33} \end{pmatrix} - A_{12} \det \begin{pmatrix} A_{21} & A_{23} \\ A_{31} & A_{33} \end{pmatrix} + A_{13} \det \begin{pmatrix} A_{21} & A_{22} \\ A_{31} & A_{32} \end{pmatrix}.
\]
Show directly that $D$ is an alternating and $3$-linear function as a function of the columns of $A$.
\exhint[Hint (Y. St\"aubli)]{Verify column linearity by expanding the $2 \times 2$ determinants and checking linearity in each of the three column vectors; verify the alternating property by checking that $D(A) = 0$ whenever two columns of $A$ coincide.}
\exinfo{This exercise is Problem 5 of Problem Sheet 16, Linear Algebra II, Spring Semester 2026.}
\exsol{sol:laplace_expansion_3x3_alternating_trilinear}
\end{exercise}

\renewcommand{\thetheorem}{20.b.\arabic{theorem}}
\setcounter{theorem}{0}
\section{Uniqueness of the Determinant}

We have shown that a determinant function exists. Now, we will prove that it is absolutely unique. Let $D : \M_{n \times n}(K) \to K$ be an $n$-linear alternating function.
Let $A \in \M_{n \times n}(K)$ be a matrix with rows $\alpha_1, \dots, \alpha_n$. We denote the identity matrix as $I$, whose rows are the standard basis vectors $\epsilon_1, \dots, \epsilon_n$ of $K^n_{\text{row}}$.

Clearly, we can express each row of $A$ in terms of the standard basis: $\alpha_i = \sum_{j=1}^n A(i,j)\epsilon_j$ for $1 \le i \le n$. Using the $n$-linearity of $D$, we can expand $D(A)$ as follows:
\begin{align*}
  D(A) &= D(\alpha_1, \dots, \alpha_n) \\
  &= D\left(\sum_{j=1}^n A(1,j)\epsilon_j, \alpha_2, \dots, \alpha_n\right) \\
  &= \sum_{j=1}^n A(1,j) \cdot D(\epsilon_j, \alpha_2, \dots, \alpha_n) \\
  &= \sum_{j=1}^n A(1,j) \cdot D\left(\epsilon_j, \sum_{k=1}^n A(2,k)\epsilon_k, \dots, \alpha_n\right) \\
  &= \sum_{j,k} A(1,j) \cdot A(2,k) \cdot D(\epsilon_j, \epsilon_k, \dots, \alpha_n)
\end{align*}
where both $j$ and $k$ run between $1$ and $n$. Continuing this expansion for all $n$ rows yields an enormous sum:
\[
D(A) = \sum_{k_1, \dots, k_n} A(1,k_1) \cdot A(2,k_2) \cdots A(n,k_n) \cdot D(\epsilon_{k_1}, \epsilon_{k_2}, \dots, \epsilon_{k_n})
\]
where each of the indices $k_1, \dots, k_n$ runs between $1$ and $n$.

Thus far, we have not used the fact that $D$ is alternating. Because $D$ is alternating, the term $D(\epsilon_{k_1}, \dots, \epsilon_{k_n}) = 0$ whenever two of the indices $k_1, \dots, k_n$ are equal. Thus, we are interested only in ordered sequences $(k_1, \dots, k_n)$ in which $k_i \in \{1, \dots, n\}$ for all $i$, and $k_i \neq k_j$ for all $i \neq j$.

Such a sequence is called a \newterm{permutation} of $\{1, \dots, n\}$, or a permutation of degree $n$. Alternatively, we can think of a permutation of degree $n$ as a bijective function $\sigma : \{1, \dots, n\} \to \{1, \dots, n\}$. Then, the sequence $(k_1, \dots, k_n) = (\sigma(1), \sigma(2), \dots, \sigma(n))$ is a permutation as defined earlier. In other words, a permutation of degree $n$ is simply a re-ordering of a list of $n$ elements: $(1, \dots, n) \mapsto (\sigma(1), \dots, \sigma(n))$.

We can nicely rewrite our expansion of $D(A)$ as:
\[
D(A) = \sum_{\sigma} A(1,\sigma(1)) \cdot A(2,\sigma(2)) \cdots A(n,\sigma(n)) \cdot D(\epsilon_{\sigma(1)}, \epsilon_{\sigma(2)}, \dots, \epsilon_{\sigma(n)})
\]
where $\sigma$ runs over all possible permutations of degree $n$.

\subsection{Properties of Permutations}

An important property about permutations is that if $\sigma$ is a permutation of degree $n$, one can pass from the sequence $(1, \dots, n)$ to $(\sigma(1), \dots, \sigma(n))$ by performing a finite sequence of interchanges of pairs.

More precisely, a \newterm{transposition} is a permutation $\theta : \{1, \dots, n\} \to \{1, \dots, n\}$ for which there exist $i, j \in \{1, \dots, n\}$ with $i \neq j$ such that $\theta(i) = j$, $\theta(j) = i$, and $\theta(k) = k$ for all $k \neq i,j$.

Every permutation $\sigma$ can be written as a composition of transpositions:
\begin{equation} \label{eq:permutation_decomp}
  \sigma = \theta_1 \circ \cdots \circ \theta_m
\end{equation}
where $\theta_1, \dots, \theta_m$ are transpositions. However, these $\theta_i$'s and the total number $m$ are not unique. The same permutation $\sigma$ can be written as $\sigma = \theta_1 \circ \cdots \circ \theta_m = \theta'_1 \circ \cdots \circ \theta'_{m'}$.

\begin{remark}
  Why is the decomposition \eqref{eq:permutation_decomp} always possible? Start with the sorted list $(1, \dots, n)$. If $\sigma(1) \neq 1$, then interchange $1$ and $\sigma(1)$. We now have $(\sigma(1), \dots, 1, \dots, n)$. Next, look at the second position. If it is not $\sigma(2)$, interchange it with $\sigma(2)$, and proceed similarly down the line. Of course, there are other ways to do it. For example, the permutation $(3,2,1)$ can be obtained in 3 steps:
  \[ (1,2,3) \to (2,1,3) \to (2,3,1) \to (3,2,1) \]
  But, it can also be obtained in just one step: $(1,2,3) \to (3,2,1)$.
\end{remark}

\begin{lemma}[Parity of Transpositions]\label{lem:permutation_parity}
  If $(\sigma(1), \dots, \sigma(n))$ is obtained from $(1, \dots, n)$ in two different ways—one time by a sequence of $m$ transpositions, and another time by a sequence of $m'$ transpositions—then $m$ and $m'$ must have the same parity. That is, $m'$ is even if and only if $m$ is even, and $m'$ is odd if and only if $m$ is odd.
\end{lemma}

In other words, $(-1)^{m'} = (-1)^m$. The parity depends only on the permutation $\sigma$, and not on the particular sequence of transpositions used.

A permutation $\sigma$ is called \newterm{even} if $m$ is even, and \newterm{odd} if $m$ is odd. The number $\sgn(\sigma) := (-1)^m \in \{-1, 1\}$ is called the \newterm{sign} of $\sigma$.

\begin{proof}[Proof of \Cref{lem:permutation_parity}]
  We know from previous results (\cref{cor:det_exists}) that determinant functions $E : \M_{n \times n}(K) \to K$ exist for $n \ge 1$. Let us fix one such determinant function $E$ (we denote it by $E$ to distinguish it from $D$ in the previous discussion).

  Consider the evaluation $E(\epsilon_{\sigma(1)}, \dots, \epsilon_{\sigma(n)})$. If $(\sigma(1), \dots, \sigma(n))$ can be obtained from $(1, \dots, n)$ by a sequence of $m$ transpositions, then the matrix with rows $\epsilon_{\sigma(1)}, \dots, \epsilon_{\sigma(n)}$ can be obtained from the identity matrix $I$ by a sequence of $m$ interchanges of pairs of rows.

  Because $E$ is alternating, swapping two rows flips the sign of the output. Therefore:
  \[
  E(\epsilon_{\sigma(1)}, \dots, \epsilon_{\sigma(n)}) = (-1)^m E(\epsilon_1, \dots, \epsilon_n) = (-1)^m E(I) = (-1)^m.
  \]
  But, this exact same logic applies to the sequence of $m'$ transpositions. Thus:
  \[
  E(\epsilon_{\sigma(1)}, \dots, \epsilon_{\sigma(n)}) = (-1)^{m'} E(\epsilon_1, \dots, \epsilon_n) = (-1)^{m'} E(I) = (-1)^{m'}.
  \]
  This identically implies that $(-1)^m = (-1)^{m'}$, proving the lemma.
\end{proof}

% Creator: Opus 5
% The order of results here looks circular at a glance and is not. Worth
% saying once, because a reader who notices it has no way to check without
% tracing the whole section.
\begin{remark}[No Circularity Here]
\label{rem:no_circularity_sign}
The argument just given proves a fact about permutations by using determinants, and the Leibniz formula below will express the determinant using $\sgn$. That is not a circle. What the proof needs is only \cref{cor:det_exists}, the \emph{existence} of some determinant function, and that was obtained from \cref{thm:recursive_det} by a recursion on $n$ in which the sign of a permutation never appears: the signs there are the explicit factors $(-1)^{i+j}$ attached to positions, not to permutations. Only after $\sgn$ has been shown to be well defined is it used, in \cref{thm:unique_det}, to write down the one determinant function in closed form.
\end{remark}

Returning to our original function $D(A)$, we have seen that:
\[
D(A) = \sum_{\sigma} A(1,\sigma(1)) \cdots A(n,\sigma(n)) \cdot D(\epsilon_{\sigma(1)}, \dots, \epsilon_{\sigma(n)}).
\]
By the alternating property of $D$, performing the transpositions necessary to sort $(\sigma(1), \dots, \sigma(n))$ back to $(1, \dots, n)$ yields:
\begin{equation} \label{eq:det_relation}
  D(\epsilon_{\sigma(1)}, \dots, \epsilon_{\sigma(n)}) = \sgn(\sigma) \cdot D(\epsilon_1, \dots, \epsilon_n) = \sgn(\sigma) \cdot D(I).
\end{equation}

\begin{theorem}[Leibniz Formula \& Uniqueness]
  \label{thm:unique_det}
  For all $n \ge 1$, there exists a unique determinant function on $\M_{n \times n}(K)$. We will denote it by $\det$. It is mathematically given by the formula:
  \[
  \det(A) = \sum_{\sigma} \sgn(\sigma) \cdot A(1,\sigma(1)) \cdot A(2,\sigma(2)) \cdots A(n,\sigma(n)).
  \]
\end{theorem}

\begin{proof}
  Existence is \cref{cor:det_exists}, so only uniqueness is at issue, and that follows immediately from \eqref{eq:det_relation}. Indeed, if $D$ is a determinant function, then $D(I) = 1$. Substituting this into \eqref{eq:det_relation} yields the exact formula for $D(A)$ stated in the theorem, proving that the function is entirely unique.
\end{proof}

\begin{corollary}[Scaling of Alternating Functions]
  \label{cor:det_pullout}
  For any $n$-linear alternating function $D : \M_{n \times n}(K) \to K$, we have $D(A) = \det(A) \cdot D(I)$ for all $A \in \M_{n \times n}(K)$.
\end{corollary}

% Creator: Opus 5
% Stated without proof, although it is the engine of three later results:
% multiplicativity, the block-matrix formula and \cref{cor:det_columns}. It is
% two lines, and the point is that the expansion above never used D(I) = 1.
\begin{proof}
  Note that the expansion leading to \eqref{eq:det_relation} used only that $D$ is $n$-linear and alternating, never that $D(I) = 1$. Substituting \eqref{eq:det_relation} into that expansion gives
  \[
  D(A) = \sum_{\sigma} A(1,\sigma(1)) \cdots A(n,\sigma(n)) \cdot \sgn(\sigma) \cdot D(I) = \left( \sum_{\sigma} \sgn(\sigma) A(1,\sigma(1)) \cdots A(n,\sigma(n)) \right) D(I),
  \]
  and the bracket is $\det(A)$ by \cref{thm:unique_det}.
\end{proof}

\begin{exercise}[Determinant of Triangular Matrices via the Leibniz Formula \faHeart]
  \label{exc:triangular_matrix_determinant}
  An $n \times n$ matrix $A$ is called \emph{upper triangular} if $A_{ij} = 0$ for all $i > j$, and \emph{lower triangular} if $A_{ij} = 0$ for all $i < j$. Show using the Leibniz formula that the determinant of any triangular matrix is given by the product of its diagonal entries:
  \[
  \det(A) = A_{11} A_{22} \cdots A_{nn}.
  \]
  \exhint[Hint (Y. St\"aubli)]{Use the Leibniz formula $\det A = \sum_{\sigma \in S_n} \sgn(\sigma) A(1,\sigma(1)) \cdots A(n,\sigma(n))$. For an upper triangular matrix ($A_{ij} = 0$ for $i > j$), show that if $\sigma \neq \id$, there must exist an index $i \in \{1, \dots, n\}$ such that $\sigma(i) < i$, which forces the corresponding factor in the product to be zero. A similar argument applies to lower triangular matrices.}
  \exinfo{This exercise is Problem 7 of Problem Sheet 17, Linear Algebra II, Spring Semester 2026, ETH Zurich (Lecturer: Prof. P. Biran).}
\exsol{sol:triangular_matrix_determinant}
\end{exercise}

\subsection{The Symmetric Group}

Denote by $S_n$ the set of permutations $\sigma : \{1, \dots, n\} \to \{1, \dots, n\}$. The set $S_n$ has the algebraic structure of a group under composition: if $\sigma, \tau \in S_n$, their composition is $(\sigma \circ \tau)(x) = \sigma(\tau(x))$. The identity function $\id$ is the neutral element. The inverse $\sigma^{-1}$ exists because $\sigma$ is bijective. Furthermore, the size of the symmetric group is $|S_n| = n!$.

% Extractor: Opus 5
\begin{exercise}[The Order of the Symmetric Group]
  \label{exc:size_of_Sn}
  Show that $|S_n| = n!$.
\exinfo{This exercise is taken from Prof.~Biran's handwritten lecture notes.}
\exsol{sol:size_of_Sn}
\end{exercise}

% Extractor: Gemini 3.1 Pro

\begin{lemma}[Multiplicativity of Signum]
  \label{lem:sgn_multiplicative}
  For all $\sigma, \tau \in S_n$, we have $\sgn(\sigma \circ \tau) = \sgn(\sigma) \cdot \sgn(\tau)$.
\end{lemma}

\begin{proof}
  Suppose $\sigma = \theta_1 \circ \cdots \circ \theta_m$ and $\tau = \lambda_1 \circ \cdots \circ \lambda_{\ell}$, where the $\theta_i$'s and $\lambda_j$'s are all transpositions. Then, their composition is simply $\sigma \circ \tau = \theta_1 \circ \cdots \circ \theta_m \circ \lambda_1 \circ \cdots \circ \lambda_{\ell}$, which explicitly consists of $m + \ell$ transpositions. Therefore:
  \[
  \sgn(\sigma \circ \tau) = (-1)^{m+\ell} = (-1)^m \cdot (-1)^{\ell} = \sgn(\sigma) \cdot \sgn(\tau).
  \]
\end{proof}

\begin{exercise}[The Symmetric Group $S_4$ and Even Permutations in $S_n$ \faHeart]
  \label{exc:symmetric_group_s4_and_alternating}
  \leavevmode
  \begin{enumerate}[label=\textbf{(\alph*)}]
    \item Explicitly list all $24$ permutations of degree $4$ in $S_4$, specify which are even and which are odd, and use this to write out the complete Leibniz formula
    \[
    \det(A) = \sum_{\sigma \in S_4} \sgn(\sigma) A(1,\sigma(1)) A(2,\sigma(2)) A(3,\sigma(3)) A(4,\sigma(4))
    \]
    for the determinant of a $4 \times 4$ matrix $A$. Notice that for $n \ge 4$, inspecting diagonal products alone is not sufficient to determine the determinant.
    \item How many even permutations are contained in $S_n$ for an arbitrary integer $n \ge 1$?
  \end{enumerate}
  \exhint[Hint (Y. St\"aubli)]{\textbf{(a)} Permutations of the same cycle structure have the same sign (e.g., $4$-cycles have sign $-1$, products of two disjoint transpositions have sign $+1$, $3$-cycles have sign $+1$, single transpositions have sign $-1$, and the identity has sign $+1$). List all $4! = 24$ permutations grouped by cycle type to write down the $24$ terms. \textbf{(b)} The case $n=1$ is trivial. For $n \ge 2$, fix a transposition $\tau \in S_n$ (e.g., $\tau = (1\;2)$) and show that the map $\sigma \mapsto \tau \circ \sigma$ is a bijection between the set of even permutations and the set of odd permutations.}
  \exinfo{This exercise is Problem 6 of Problem Sheet 17, Linear Algebra II, Spring Semester 2026, ETH Zurich (Lecturer: Prof. P. Biran).}
\exsol{sol:symmetric_group_s4_and_alternating}
\end{exercise}

\subsection{Multiplicativity of the Determinant}

\begin{theorem}[Multiplicativity of the Determinant]
  \label{thm:det_multiplicative}
  For any matrices $A, B \in M_{n \times n}(K)$, we have $\det(A \cdot B) = \det(A) \cdot \det(B)$.
\end{theorem}

\begin{proof}
  Fix a matrix $B \in \M_{n \times n}(K)$, and consider a function $D : \M_{n \times n}(K) \to K$ defined by $D(A) := \det(A \cdot B)$.

  Claim: $D$ is $n$-linear and alternating.

  Proof of claim: Let the rows of $A$ be $\alpha_1, \dots, \alpha_n$. If we view $\alpha_i$ as a $1 \times n$ matrix, then the matrix product $\alpha_i \cdot B$ is also $1 \times n$. In fact, the rows of the product matrix $A \cdot B$ are exactly $\alpha_1 \cdot B, \dots, \alpha_n \cdot B$. So, $D(A) = \det(\alpha_1 \cdot B, \dots, \alpha_n \cdot B)$.

  Clearly, for two row vectors $\alpha_i$ and $\alpha'_i$, we have $(\alpha_i + \alpha'_i) \cdot B = \alpha_i \cdot B + \alpha'_i \cdot B$. Also, for any scalar $c \in K$, $(c\alpha_i) \cdot B = c(\alpha_i \cdot B)$. Since the determinant function $\det$ is $n$-linear with respect to its rows, it heavily follows that $D$ is also $n$-linear.

  Furthermore, if two rows of $A$ are equal (i.e., $\alpha_i = \alpha_j$), then the corresponding rows of $A \cdot B$ are also equal ($\alpha_i \cdot B = \alpha_j \cdot B$). Because $\det$ is alternating, $D(A) = \det(\alpha_1 \cdot B, \dots, \alpha_n \cdot B) = 0$. This rigorously proves the claim that $D$ is alternating.

  We continue with the proof of \cref{thm:det_multiplicative}. By \Cref{cor:det_pullout}, any $n$-linear alternating function satisfies $D(A) = \det(A) \cdot D(I)$. We compute $D(I)$ using our definition: $D(I) = \det(I \cdot B) = \det(B)$. Substituting this back yields $D(A) = \det(A) \cdot \det(B)$, completing the proof.
\end{proof}

\begin{corollary}[Determinant of an Inverse]
  \label{cor:det_inverse}
  If $A$ is invertible, then $\det(A) \neq 0$. Moreover, $\det(A^{-1}) = \frac{1}{\det(A)}$.
\end{corollary}

\begin{proof}
  Assume $A$ is invertible, and denote its inverse by $A^{-1}$. Since $A \cdot A^{-1} = I$, taking the determinant of both sides and applying \cref{thm:det_multiplicative} yields $\det(A) \cdot \det(A^{-1}) = \det(I) = 1$. This automatically implies that $\det(A)$ cannot be zero, and dividing by $\det(A)$ yields the expected formula.
\end{proof}

\begin{exercise}[Direct Proof of Multiplicativity of the Determinant via the Leibniz Formula]
  \label{exc:multiplicativity_via_leibniz}
  Prove the following proposition directly using the Leibniz formula:
  \begin{proposition*}
    Let $K$ be a field and let $A, B \in \M_{n \times n}(K)$. Then
    \[
    \det(A B) = \det(A) \cdot \det(B).
    \]
  \end{proposition*}
  \exhint[Official Hint]{Let $\{e_i\}_{i=1}^n$ denote the standard basis of $K^n$ and write the matrix $B$ as a list of column vectors $B = \left( \sum_{s_1=1}^n B(s_1, 1) e_{s_1}, \dots, \sum_{s_n=1}^n B(s_n, n) e_{s_n} \right)$. Expand $\det(AB) = \det(A B e_1, \dots, A B e_n)$ using the multilinearity of the determinant:
  \[
  \det(AB) = \sum_{s_1=1}^n \cdots \sum_{s_n=1}^n \left( \prod_{i=1}^n B(s_i, i) \right) \det(A e_{s_1}, \dots, A e_{s_n}).
  \]
  Show that terms where the indices $s_1, \dots, s_n$ are not distinct vanish, so the sum reduces to permutations $\sigma \in S_n$. Then prove and apply the lemma:
  \begin{lemma*}
    For all $A \in \M_{n \times n}(K)$ and all $\sigma \in S_n$,
    \[
    \det(A e_{\sigma(1)}, \dots, A e_{\sigma(n)}) = \sgn(\sigma) \det(A).
    \]
  \end{lemma*}}
  \exinfo{This exercise is Problem 10 of Problem Sheet 17, Linear Algebra II, Spring Semester 2026, ETH Zurich (Lecturer: Prof. P. Biran).}
\exsol{sol:multiplicativity_via_leibniz}
\end{exercise}

\begin{exercise}[General Properties of Determinants (Multiple Choice)]
  \label{exc:general_determinant_properties_mc}
  Which of the following statements are correct for arbitrary matrices $A, B \in \M_{n \times n}(\mathbb{R})$ with $n \ge 2$?
  \begin{enumerate}[label=\textbf{(\alph*)}]
    \item $\det(A B) = \det(A) \det(B)$.
    \item If $\det(A) \neq 0$, then the column vectors $a_1, \dots, a_n$ of $A$ are linearly independent.
    \item $\det(A B) = \det(B A)$.
    \item For every non-zero real number $\lambda$, $\det(\lambda A) = \lambda \det(A)$.
    \item $\det(A + B) = \det(A) + \det(B)$.
  \end{enumerate}
  \exinfo{This exercise is Multiple Choice Question 1 of Problem Sheet 17, Linear Algebra II, Spring Semester 2026, ETH Zurich (Lecturer: Prof. P. Biran).}
\exsol{sol:general_determinant_properties_mc}
\end{exercise}

\begin{exercise}[Properties of the Determinant Mapping (Single Choice)]
  \label{exc:determinant_properties_linearity_sc}
  Let $K$ be a field and let $\M_{n \times n}(K)$ be the vector space of $n \times n$ matrices over $K$. Which of the following statements is in general \textbf{false}?
  \begin{enumerate}[label=\textbf{(\alph*)}]
    \item A square matrix $A$ over $K$ is invertible if and only if $\det(A) \neq 0$.
    \item The determinant of an upper triangular matrix depends only on its diagonal entries.
    \item For every $n \ge 0$, the determinant $\det : \M_{n \times n}(K) \to K$ is a linear map.
    \item For every $n > 0$, the determinant mapping $\det : \M_{n \times n}(K) \to K$ is surjective.
  \end{enumerate}
  \exinfo{This exercise is Single Choice Question 2 of Problem Sheet 17, Linear Algebra II, Spring Semester 2026, ETH Zurich (Lecturer: Prof. P. Biran).}
\exsol{sol:determinant_properties_linearity_sc}
\end{exercise}

\renewcommand{\thetheorem}{20.c.\arabic{theorem}}
\setcounter{theorem}{0}
\section{More on Determinants: Properties of Determinants and the Transpose}

Recall that for any matrix $M \in M_{m \times n}(K)$, we have the transposed matrix $M\transp \in M_{n \times m}(K)$, where the entries are given by $M\transp(i,j) = M(j,i)$.

\begin{theorem}[Determinant of the Transpose]
  \label{thm:det_transpose}
  For all $A \in \M_{n \times n}(K)$, we have $\det(A\transp) = \det(A)$.
\end{theorem}

\begin{proof}
  If $\sigma$ is a permutation of degree $n$, then by definition of the transpose, $A\transp(i, \sigma(i)) = A(\sigma(i), i)$.
  Therefore, expanding the determinant definition gives:
  \[
  \det(A\transp) = \sum_{\sigma \in S_n} \sgn(\sigma) \cdot A(\sigma(1), 1) \cdots A(\sigma(n), n).
  \]
  If we let $j = \sigma(i)$, then $i = \sigma^{-1}(j)$, and consequently $A(\sigma(i), i) = A(j, \sigma^{-1}(j))$. Hence, we can rearrange the product completely:
  \begin{equation} \label{eq:transpose_prod}
    A(\sigma(1), 1) \cdots A(\sigma(n), n) = A(1, \sigma^{-1}(1)) \cdots A(n, \sigma^{-1}(n))
  \end{equation}
  This holds because every factor on the left-hand side of \eqref{eq:transpose_prod} will appear exactly once on the right-hand side, and vice-versa.

  Now, since $\sigma \circ \sigma^{-1} = \id$, we have $\sgn(\sigma) \cdot \sgn(\sigma^{-1}) = \sgn(\id) = 1$, which implies $\sgn(\sigma) = \sgn(\sigma^{-1})$.
  When $\sigma$ runs over $S_n$ (all possible permutations of degree $n$), its inverse $\sigma^{-1}$ also naturally runs over the entirety of $S_n$.

  This implies:
  \begin{align*}
    \det(A\transp) &= \sum_{\sigma \in S_n} \sgn(\sigma) A(\sigma(1), 1) \cdots A(\sigma(n), n) \\
    &= \sum_{\sigma \in S_n} \sgn(\sigma^{-1}) A(1, \sigma^{-1}(1)) \cdots A(n, \sigma^{-1}(n)) \\
    &= \sum_{\tau \in S_n} \sgn(\tau) A(1, \tau(1)) \cdots A(n, \tau(n)) \\
    &= \det(A).
  \end{align*}
\end{proof}

\begin{example}[Determinant of a Transposed Matrix]
  Let $A = \left(\begin{smallmatrix} 1 & 2 \\ 3 & 4 \end{smallmatrix}\right)$. Its transpose is $A\transp = \left(\begin{smallmatrix} 1 & 3 \\ 2 & 4 \end{smallmatrix}\right)$. We have $\det(A) = (1)(4) - (2)(3) = -2$, and $\det(A\transp) = (1)(4) - (3)(2) = -2$. Thus, the theorem holds.
\end{example}

\begin{corollary}[Column Alternating Property]
  \label{cor:det_columns}
  If $D : \M_{n \times n}(K) \to K$ is an alternating $n$-linear function of the rows of the matrix, then $D$ is also an alternating $n$-linear function of the columns of the matrix. In fact, $D(A\transp) = D(A)$ for all $A \in \M_{n \times n}(K)$. In particular, this holds for $D = \det$.
\end{corollary}

\begin{proof}
  For any $A \in M_{n \times n}(K)$, we have previously established $D(A) = \det(A) \cdot D(I)$. Applying this to $A\transp$ gives $D(A\transp) = \det(A\transp) \cdot D(I)$. Since $\det(A\transp) = \det(A)$, this evaluates exactly to $\det(A) \cdot D(I) = D(A)$.
\end{proof}

Because the determinant acts symmetrically on rows and columns, we can perform column operations just as we do row operations.

\begin{theorem}[Invariance Under Row Operations]
  \label{thm:row_operations}
  Let $A \in \M_{n \times n}(K)$.
  \begin{enumerate}[label=\textbf{(\alph*)}]
    \item If $B$ is obtained from $A$ by adding a multiple of one row of $A$ to another row (i.e., $R_i + c \cdot R_j \to R_i$ for $i \neq j$), or by doing the analogous operation on columns, then $\det(B) = \det(A)$.
    \item If $B$ is obtained from $A$ by interchanging two rows (i.e., $R_i \leftrightarrow R_j$ for $i \neq j$), or columns, then $\det(B) = -\det(A)$.
    \item If $B$ is obtained from $A$ by multiplying a row, or a column, by a scalar (i.e., $c \cdot R_i \to R_i$), then $\det(B) = c \det(A)$.
  \end{enumerate}
\end{theorem}

\begin{proof}
  Statement \textbf{(b)} follows directly from the fact that the determinant is an alternating function. Statement \textbf{(c)} follows from the $n$-linearity of the determinant.

  To prove \textbf{(a)}, let $A$ have rows $\alpha_1, \dots, \alpha_n$. Assume $j < i$. Expanding by linearity on the $i$-th row gives:
  \begin{align*}
    \det(B) &= \det(\alpha_1, \dots, \alpha_i + c\alpha_j, \dots, \alpha_j, \dots, \alpha_n) \\
    &= \det(\alpha_1, \dots, \alpha_i, \dots, \alpha_n) + c \det(\alpha_1, \dots, \alpha_j, \dots, \alpha_j, \dots, \alpha_n) \\
    &= \det(A) + 0 = \det(A).
  \end{align*}
  The second term identically vanishes because two rows are equal ($\alpha_j$ appears twice). If $j > i$, the proof is similar.

  In each of the three statements, the assertion about columns follows from the one about rows: a column operation on $A$ is the corresponding row operation on $A\transp$, and $\det(A\transp) = \det(A)$ by \cref{thm:det_transpose}.
\end{proof}

\begin{exercise}[Row Operations and Determinants \faHeart]
  \label{exc:row_operations_and_determinant_divisibility}
  \leavevmode
  \begin{enumerate}[label=\textbf{(\alph*)}]
    \item Let $K$ be a field, $\lambda \in K$, and $A \in \M_{n \times n}(K)$. Show:
    \begin{enumerate}[label=\textbf{(\roman*)}]
      \item If $B$ is obtained from $A$ by multiplying row $i$ by $\lambda$ ($\lambda \cdot R_i \to R_i$), then $\det B = \lambda \det A$.
      \item If $B$ is obtained from $A$ by swapping rows $i$ and $j$ ($R_i \leftrightarrow R_j$), then $\det B = -\det A$.
      \item If $B$ is obtained from $A$ by adding $\lambda$ times row $i$ to row $j$ ($R_j + \lambda \cdot R_i \to R_j$ with $i \neq j$), then $\det B = \det A$.
    \end{enumerate}
    \item The numbers $2014, 1484, 3710$, and $6996$ are all divisible by $106$. Show without explicitly computing the numerical value that
    \[
    \det \begin{pmatrix}
      2 & 1 & 3 & 6 \\
      0 & 4 & 7 & 9 \\
      1 & 8 & 1 & 9 \\
      4 & 4 & 0 & 6
    \end{pmatrix}
    \]
    is also divisible by $106$.
  \end{enumerate}
  \exhint[Hint (Y. St\"aubli)]{\textbf{(a)} These are the fundamental computational rules for determinants under elementary row operations; see \cref{thm:row_operations}. \textbf{(b)} Perform the elementary row operations $1000 R_1 + R_4 \to R_4$, $100 R_2 + R_4 \to R_4$ and $10 R_3 + R_4 \to R_4$ to obtain the fourth row $(2014, 1484, 3710, 6996)$, and observe that a common factor in any row can be factored out of the determinant.}
  \exinfo{This exercise is Problem 2 of Problem Sheet 17, Linear Algebra II, Spring Semester 2026, ETH Zurich (Lecturer: Prof. P. Biran).}
\exsol{sol:row_operations_and_determinant_divisibility}
\end{exercise}

\begin{exercise}[Invariance Under Matrix Operations (Single Choice)]
  \label{exc:invariance_under_matrix_operations_sc}
  Under which of the following operations does the determinant of a matrix in general \textbf{not} remain unchanged?
  \begin{enumerate}[label=\textbf{(\alph*)}]
    \item Interchanging two rows.
    \item Adding a scalar multiple of one row to another row.
    \item Transposing the matrix.
    \item Replacing the matrix by a similar matrix.
  \end{enumerate}
  \exinfo{This exercise is Single Choice Question 3 of Problem Sheet 17, Linear Algebra II, Spring Semester 2026, ETH Zurich (Lecturer: Prof. P. Biran).}
\exsol{sol:invariance_under_matrix_operations_sc}
\end{exercise}

\begin{exercise}[Row Transformations of a $2 \times 2$ Determinant (Multiple Choice)]
  \label{exc:row_transformations_2x2_mc}
  Suppose $\det \begin{pmatrix} a & b \\ c & d \end{pmatrix} = 4$. Which of the following statements are correct?
  \begin{enumerate}[label=\textbf{(\alph*)}]
    \item $\det \begin{pmatrix} 2a & 2b \\ 2c & 2d \end{pmatrix} = 8$.
    \item $\det \begin{pmatrix} a & b \\ c-a & d-b \end{pmatrix} = 4$.
    \item $\det \begin{pmatrix} a & b \\ c+2a & d+2b \end{pmatrix} = 4$.
    \item $\det \begin{pmatrix} a & b \\ 3c & 3d \end{pmatrix} = 12$.
  \end{enumerate}
  \exinfo{This exercise is Multiple Choice Question 2 of Problem Sheet 17, Linear Algebra II, Spring Semester 2026, ETH Zurich (Lecturer: Prof. P. Biran).}
\exsol{sol:row_transformations_2x2_mc}
\end{exercise}

\subsection{Determinants of Block Matrices}

Calculating the determinant of a large matrix by expanding all permutations is computationally expensive. Block matrices provide a powerful shortcut. Consider a block matrix $M$ of the form $M = \left(\begin{smallmatrix} A & B \\ 0 & C \end{smallmatrix}\right)$, where $A \in \M_{r \times r}(K)$, $C \in \M_{s \times s}(K)$, and $B \in \M_{r \times s}(K)$.

\begin{proposition}[Determinant of Block Matrices]
  \label{prop:block_matrix}
  For a block matrix $M$ defined as above, $\det(M) = \det(A) \cdot \det(C)$, where $A \in \M_{r \times r}(K)$, $C \in \M_{s \times s}(K)$, and $B \in \M_{r \times s}(K)$.
\end{proposition}

\begin{proof}
  Define a helper function $D(A, B, C) := \det \left(\begin{smallmatrix} A & B \\ 0 & C \end{smallmatrix}\right)$ as a function $D : \M_{r \times r}(K) \times \M_{r \times s}(K) \times \M_{s \times s}(K) \to K$.

  If we fix the first two matrices $A$ and $B$, then the function mapping $C \mapsto D(A, B, C)$ is $s$-linear and alternating. By \Cref{cor:det_pullout}, any such function satisfies:
  \begin{equation} \label{eq:block_step1}
    D(A, B, C) = \det(C) \cdot D(A, B, I_s).
  \end{equation}

  Let us calculate $D(A, B, I_s) = \det \left ( \begin{smallmatrix} A & B \\ 0 & I_s \end{smallmatrix}\right)$. By doing row operations of the type $R_i + c R_j \to R_i$ (where $j \neq i$), we can use the last $s$ rows of the identity matrix $I_s$ to eliminate the entries in $B$ entirely. This transforms the matrix into $\left(\begin{smallmatrix} A & 0 \\ 0 & I_s \end{smallmatrix}\right)$.

  Since these operations do not change the determinant, we have:
  \begin{equation} \label{eq:block_step2}
    \det \begin{pmatrix} A & B \\ 0 & I_s \end{pmatrix} = \det \begin{pmatrix} A & 0 \\ 0 & I_s \end{pmatrix}.
  \end{equation}

  The function $A \mapsto \det \left(\begin{smallmatrix} A & 0 \\ 0 & I_s \end{smallmatrix}\right)$ is $r$-linear and alternating. By \cref{cor:det_pullout} again, it evaluates to $\det(A) \cdot \det \left(\begin{smallmatrix} I_r & 0 \\ 0 & I_s \end{smallmatrix}\right) = \det(A) \cdot \det(I_{r+s}) = \det(A) \cdot 1 = \det(A)$.

  Going back through \eqref{eq:block_step2} and then to \eqref{eq:block_step1}, we finally get that $D(A, B, C) = \det(A) \cdot \det(C)$.
\end{proof}

\begin{exercise}[The Other Block Shape]
  \label{exc:block_lower}
  Show that $\det \left(\begin{smallmatrix} A & 0 \\ B & C \end{smallmatrix}\right) = \det(A) \cdot \det(C)$, where $A \in \M_{r \times r}(K)$, $C \in \M_{s \times s}(K)$, and $B \in \M_{s \times r}(K)$. (Hint: Use column operations to reduce $B$ to zero, or consider the transpose of the matrix and apply \cref{prop:block_matrix}.)
\exinfo{This exercise is taken from Prof.~Biran's handwritten lecture notes.}
\exsol{sol:block_lower}
\end{exercise}

\begin{exercise}[Determinant of a Block Matrix with Commuting Blocks]
  \label{exc:block_matrix_commuting_blocks}
  Let $K$ be a field and let $A, B, C, D \in \M_{n \times n}(K)$. Assume that $A$ and $C$ commute ($A C = C A$) and that $\det A \neq 0$. Show that
  \[
  \det \begin{pmatrix} A & B \\ C & D \end{pmatrix} = \det(A D - C B).
  \]
  \exhint[Official Hint]{Multiply the given block matrix from the left by $\begin{pmatrix} I_n & 0 \\ -C & A \end{pmatrix}$. Use the block matrix determinant formula $\det \begin{pmatrix} X & 0 \\ Y & Z \end{pmatrix} = \det(X) \det(Z)$ and the multiplicativity of the determinant, together with the commutativity relation $A C = C A$, to conclude.}
  \exinfo{This exercise is Problem 9 of Problem Sheet 17, Linear Algebra II, Spring Semester 2026, ETH Zurich (Lecturer: Prof. P. Biran).}
\exsol{sol:block_matrix_commuting_blocks}
\end{exercise}

\begin{example}[Calculation of a $4 \times 4$ Determinant]
  Let $A = \left(\begin{smallmatrix} 1 & -1 & 2 & 3 \\ 2 & 2 & 0 & 2 \\ 4 & 1 & -1 & -1 \\ 1 & 2 & 3 & 0 \end{smallmatrix}\right) \in \M_{4 \times 4}(\mathbb{R})$.
  We begin by performing row operations $R_2 - 2R_1 \to R_2$, $R_3 - 4R_1 \to R_3$, and $R_4 - R_1 \to R_4$. This transforms $A$ into:
  \[
  \begin{pmatrix}
    1 & -1 & 2 & 3 \\
    0 & 4 & -4 & -4 \\
    0 & 5 & -9 & -13 \\
    0 & 3 & 1 & -3
  \end{pmatrix}
  \]
  Now, we clean the second column below the diagonal using $R_3 - \frac{5}{4}R_2 \to R_3$ and $R_4 - \frac{3}{4}R_2 \to R_4$. This reveals a beautiful block structure:
  \[
  \left(
  \begin{array}{cc|cc}
    1 & -1 & 2 & 3 \\
    0 & 4 & -4 & -4 \\ \hline
    0 & 0 & -4 & -8 \\
    0 & 0 & 4 & 0
  \end{array}
  \right)
  \]

  By applying \cref{prop:block_matrix} to this partitioned form, we can see that the determinant of the $4 \times 4$ matrix is simply the product of the determinants of the $2 \times 2$ diagonal blocks:
  \[
  \det(A) = \det \begin{pmatrix} 1 & -1 \\ 0 & 4 \end{pmatrix} \cdot \det \begin{pmatrix} -4 & -8 \\ 4 & 0 \end{pmatrix}
  \]
  Consequently, we evaluate the individual pieces:
  \[
  \det(A) = (4) \cdot (32) = 128.
  \]
\end{example}

\begin{exercise}[Determinant Computation via Gaussian Elimination]
  \label{exc:determinants_via_gaussian_elimination}
  Compute the determinants of the following matrices:
  \[
  A = \begin{pmatrix}
    0 & 1 & 0 & 0 \\
    -1 & 0 & 1 & 0 \\
    1 & 2 & -3 & 1 \\
    0 & -4 & 2 & -1
  \end{pmatrix}, \quad
  B = \begin{pmatrix}
    0 & 0 & 0 & 0 & 1 \\
    -1 & 2 & 0 & 0 & 1 \\
    0 & 1 & 0 & -1 & 0 \\
    0 & 0 & -1 & 0 & 1 \\
    0 & 0 & 0 & 1 & 0
  \end{pmatrix},
  \]
  \[
  C = \begin{pmatrix}
    2 & -3 & 5 & 1 & 4 \\
    2 & -3 & 1 & -6 & 18 \\
    4 & -3 & 9 & 6 & 10 \\
    -2 & 4 & -6 & -1 & -1 \\
    -6 & 11 & -23 & -14 & 9
  \end{pmatrix}.
  \]
  \exhint[Hint (Y. St\"aubli)]{Use elementary row operations (Gaussian elimination) to reduce each matrix to triangular form, keeping track of any row swaps and row scalings.}
  \exinfo{This exercise is Problem 3 of Problem Sheet 17, Linear Algebra II, Spring Semester 2026, ETH Zurich (Lecturer: Prof. P. Biran).}
\exsol{sol:determinants_via_gaussian_elimination}
\end{exercise}

\begin{exercise}[The Vandermonde Determinant \faHeart]
  \label{exc:vandermonde_determinant}
  Let $n \ge 2$ and let $x_1, \dots, x_n \in K$. Show that
  \[
  \det \begin{pmatrix}
    1 & x_1 & x_1^2 & \cdots & x_1^{n-1} \\
    1 & x_2 & x_2^2 & \cdots & x_2^{n-1} \\
    \vdots & \vdots & \vdots & \ddots & \vdots \\
    1 & x_n & x_n^2 & \cdots & x_n^{n-1}
  \end{pmatrix}
  = \prod_{1 \le i < j \le n} (x_j - x_i).
  \]
  \textit{Remark.} Products of this form are called \emph{Vandermonde determinants}, and the matrix is known as the \emph{Vandermonde matrix}.
  \exhint[Hint (Y. St\"aubli)]{Use induction on $n$. Subtract the first row from each of the subsequent rows, factor out $(x_i - x_1)$ from the $i$-th row for each $i \ge 2$, and use the identity $x_i^m - x_1^m = (x_i - x_1) \sum_{k=0}^{m-1} x_i^k x_1^{m-1-k}$ along with column operations to reduce the remaining $(n-1) \times (n-1)$ determinant to the Vandermonde determinant of order $n-1$.}
  \exinfo{This exercise is Problem 8 of Problem Sheet 17, Linear Algebra II, Spring Semester 2026, ETH Zurich (Lecturer: Prof. P. Biran).}
\exsol{sol:vandermonde_determinant}
\end{exercise}

\subsection{Cofactor Expansion and the Classical Adjoint}

Recall that if $D$ is a determinant function on $\M_{(n-1) \times (n-1)}(K)$, then for any fixed $j \in \{1, \dots, n\}$, the function $E_j(A) := \sum_{i=1}^n (-1)^{i+j} A_{ij} D(A(i|j))$ is a determinant function as well. But, we have seen that on $\M_{n \times n}(K)$ there is a strictly unique determinant function, $\det$.

\begin{corollary}[Laplace Expansion]
  \label{cor:cofactor_expansion}
  For all $A \in \M_{n \times n}(K)$, we have the following $n$ different formulas for $\det(A)$ (expanding along column $j$):
  \[
  \det(A) = \sum_{i=1}^n (-1)^{i+j} A_{ij} \det A(i|j), \quad \text{for any } j = 1, \dots, n.
  \]
\end{corollary}

\begin{exercise}[Determinant of a $5 \times 5$ Matrix via Laplace Expansion]
  \label{exc:determinant_5x5_laplace}
  Compute the determinant of the matrix
  \[
  B = \begin{pmatrix}
    1 & 1 & 1 & 1 & 1 \\
    0 & 0 & 0 & 2 & 3 \\
    0 & 0 & 2 & 3 & 0 \\
    0 & 2 & 3 & 0 & 0 \\
    2 & 3 & 0 & 0 & 0
  \end{pmatrix}
  \]
  over the field $\mathbb{R}$ and over the finite field $\mathbb{F}_5$. Is $B$ invertible over each of these fields?
  \exhint[Hint (Y. St\"aubli)]{Compute the determinant using Laplace expansion, for instance along rows or columns containing many zero entries.}
  \exinfo{This exercise is Problem 1 of Problem Sheet 17, Linear Algebra II, Spring Semester 2026, ETH Zurich (Lecturer: Prof. P. Biran).}
\exsol{sol:determinant_5x5_laplace}
\end{exercise}

\begin{exercise}[Determinant of a Tridiagonal Toeplitz Matrix \faHeart]
  \label{exc:tridiagonal_toeplitz_determinant}
  Let $A_n \in \M_{n \times n}(\mathbb{R})$ be the tridiagonal matrix
  \[
  A_n = \begin{pmatrix}
    2 & 1 & 0 & \cdots & 0 \\
    1 & 2 & 1 & \cdots & 0 \\
    0 & 1 & 2 & \cdots & 0 \\
    \vdots & \vdots & \ddots & \ddots & \vdots \\
    0 & 0 & \cdots & 1 & 2
  \end{pmatrix}.
  \]
  Prove that $\det(A_n) = n+1$ for all $n \ge 1$.
  \exhint[Hint (Y. St\"aubli)]{Use induction on $n$. Expand the determinant along the first row (or column) to establish the recurrence relation $\det(A_n) = 2 \det(A_{n-1}) - \det(A_{n-2})$. Since this relates $\det(A_n)$ to two preceding terms, verify two base cases ($n=1$ and $n=2$).}
  \exinfo{This exercise is Problem 4 of Problem Sheet 17, Linear Algebra II, Spring Semester 2026, ETH Zurich (Lecturer: Prof. P. Biran).}
\exsol{sol:tridiagonal_toeplitz_determinant}
\end{exercise}

\begin{exercise}[Parameter for Determinant Value (Single Choice)]
  \label{exc:parameter_for_3x3_determinant_sc}
  For which $x \in \mathbb{R}$ does the identity
  \[
  \det \begin{pmatrix} 1 & 1 & 1 \\ 1 & x & 1 \\ 1 & 1 & x \end{pmatrix} = 1
  \]
  hold?
  \begin{enumerate}[label=\textbf{(\alph*)}]
    \item $x = -2$
    \item $x = 2$
    \item $x = -1$
    \item $x = 1$
  \end{enumerate}
  \exinfo{This exercise is Single Choice Question 1 of Problem Sheet 17, Linear Algebra II, Spring Semester 2026, ETH Zurich (Lecturer: Prof. P. Biran).}
\exsol{sol:parameter_for_3x3_determinant_sc}
\end{exercise}

The scalars $C_{ij} := (-1)^{i+j} \det A(i|j)$ are called the \newterm{$i,j$ cofactor} of $A$. We can succinctly write the expansion of $\det(A)$ by the cofactors of the $j$-th column as $\det(A) = \sum_{i=1}^n A_{ij} \cdot C_{ij}$.

\begin{lemma*}[Orthogonality of Cofactors]\label{lem:false_cofactor}
  If we replace $A_{ij}$ in the formula by $A_{ik}$ (where $k$ is fixed), we always get $0$. In other words, for all $A \in \M_{n \times n}(K)$ and for $j \neq k$, we have $\sum_{i=1}^n A_{ik} \cdot C_{ij} = 0$.
\end{lemma*}

\begin{proof}
  Let $B \in \M_{n \times n}(K)$ be the matrix obtained from $A$ by replacing (not exchanging!) column $j$ of $A$ by column $k$ of $A$. So, in $B$, column $j$ equals column $k$. Clearly, $\det(B) = 0$ since $B$ has two equal columns.

  Expanding $\det(B)$ along the $j$-th column gives:
  \[
  0 = \det(B) = \sum_{i=1}^n (-1)^{i+j} B_{ij} \det B(i|j).
  \]
  By construction, $B_{ij} = A_{ik}$ for all $i$. Furthermore, since we only modified column $j$, deleting column $j$ from both matrices leaves identical submatrices: $B(i|j) = A(i|j)$ for all $i$. Thus:
  \[
  0 = \sum_{i=1}^n (-1)^{i+j} A_{ik} \det A(i|j) = \sum_{i=1}^n A_{ik} \cdot C_{ij}.
  \]
\end{proof}

\begin{summary}
\label{sum:cofactor_orthogonality}
For $A \in \M_{n \times n}(K)$, we have $\sum_{i=1}^n A_{ik} \cdot C_{ij} = \delta_{jk} \cdot \det(A)$, where $\delta_{jk}$ is the Kronecker delta (which is 1 if $j=k$, and 0 otherwise).
\end{summary}

\begin{definition*}[Classical Adjoint]\label{def:classical_adjoint}
  The transpose of the matrix of cofactors, $(C_{ij})\transp$, is called the \newterm{classical adjoint} (or \newterm{adjugate}) of $A$, denoted $\adj A \in \M_{n \times n}(K)$. Explicitly, $(\adj A)_{ij} = C_{ji} = (-1)^{i+j} \det A(j|i)$.
\end{definition*}

\Cref{sum:cofactor_orthogonality} immediately implies the matrix equation:
\begin{equation} \label{eq:adj_property_right}
  (\adj A) \cdot A = (\det A) \cdot I
\end{equation}
The bookkeeping is worth doing once. The entry of $(\adj A) \cdot A$ in place $(j,k)$ is
\[
\big( (\adj A) \cdot A \big)_{jk} = \sum_{i=1}^{n} (\adj A)_{ji} \, A_{ik} = \sum_{i=1}^{n} C_{ij} \, A_{ik} = \delta_{jk} \cdot \det A,
\]
using $(\adj A)_{ji} = C_{ij}$ from the Definition on \cpageref{def:classical_adjoint} for the middle equality and \cref{sum:cofactor_orthogonality} for the last, and $\delta_{jk} \det A$ is exactly the entry of $(\det A) \cdot I$ in that place.

\begin{lemma*}[Left Adjugate Identity]\label{lem:adj_property_left}
  It is also true that $A \cdot (\adj A) = (\det A) \cdot I$.
\end{lemma*}

\begin{proof}
  Since $A\transp(i|j) = (A(j|i))\transp$, we have
  \[(\adj A\transp)_{ij} = (-1)^{i+j} \det A\transp(j|i) = (-1)^{i+j} \det(A(i|j)\transp) = (-1)^{i+j} \det A(i|j).\]
  This perfectly matches $(\adj A)_{ji}$, which is exactly $((\adj A)\transp)_{ij}$. Thus, $\adj A\transp = (\adj A)\transp$.
  Applying \eqref{eq:adj_property_right} to the matrix $A\transp$ gives $(\adj A\transp) \cdot A\transp = (\det A\transp) \cdot I$. Using our transpose identities, this becomes $(\adj A)\transp \cdot A\transp = (\det A) \cdot I$. Taking the transpose of both sides finally yields $A \cdot (\adj A) = (\det A) \cdot I$.
\end{proof}

\begin{exercise}[Properties of the $2 \times 2$ Adjugate Matrix]
  \label{exc:properties_2x2_adjugate}
  Let $K$ be a field. For a $2 \times 2$ matrix $A = \begin{pmatrix} A_{11} & A_{12} \\ A_{21} & A_{22} \end{pmatrix} \in \M_{2 \times 2}(K)$, the adjugate of $A$ is defined by
  \[
  \adj A = \begin{pmatrix} A_{22} & -A_{12} \\ -A_{21} & A_{11} \end{pmatrix}.
  \]
  Let $\det$ be the unique determinant function on $\M_{2 \times 2}(K)$. Show:
  \begin{enumerate}[label=\textbf{(\alph*)}]
    \item $(\adj A) A = A (\adj A) = (\det A) I_2$;
    \item $\det(\adj A) = \det(A)$;
    \item $\adj(A\transp) = (\adj A)\transp$.
  \end{enumerate}
  \exhint[Hint (Y. St\"aubli)]{Perform direct algebraic calculations using the explicit entries of $A$ and $\adj A$.}
  \exinfo{This exercise is Problem 5 of Problem Sheet 17, Linear Algebra II, Spring Semester 2026, ETH Zurich (Lecturer: Prof. P. Biran).}
\exsol{sol:properties_2x2_adjugate}
\end{exercise}

\begin{theorem}[Inverse via Adjugate]
  \label{thm:inverse_adjoint}
  Let $A \in \M_{n \times n}(K)$. Then, $A$ is invertible if and only if $\det A \neq 0$. When $A$ is invertible, its inverse is given by $A^{-1} = \frac{1}{\det A} \cdot (\adj A)$.
\end{theorem}

% Creator: Opus 5
\begin{proof}
One direction is \cref{cor:det_inverse}: if $A$ is invertible then $\det A \neq 0$.

Conversely, assume $\det A \neq 0$. Then the scalar $\frac{1}{\det A} \in K$ exists, and the two adjugate identities \eqref{eq:adj_property_right} and \cpageref{lem:adj_property_left} give
\[
A \cdot \left( \frac{1}{\det A} \adj A \right) = \frac{1}{\det A} \big( A \cdot \adj A \big) = \frac{1}{\det A} (\det A) \cdot I = I,
\]
and in the same way $\left( \frac{1}{\det A} \adj A \right) \cdot A = I$. So, $\frac{1}{\det A} \adj A$ is a two-sided inverse of $A$, which makes $A$ invertible with $A^{-1} = \frac{1}{\det A} \adj A$.
\end{proof}

% Extractor: Gemini 3.1 Pro

\subsection{Cramer's Rule}

Let $A \in \M_{n \times n}(K)$. We would like an explicit formula to solve the system of linear equations $A \cdot x = b$, where $x = \left(\begin{smallmatrix} x_1 \\ \vdots \\ x_n \end{smallmatrix}\right)$ are the unknowns and $b = \left(\begin{smallmatrix} b_1 \\ \vdots \\ b_n \end{smallmatrix}\right) \in K^n$ is a given vector.

Assume $A \cdot x = b$. Multiplying both sides by the adjoint gives $(\adj A) \cdot A \cdot x = (\adj A) \cdot b$. This directly implies $(\det A) \cdot x = (\adj A) \cdot b$.
Looking at the $j$-th row of this vector equation, we get:
\begin{align*}
  (\det A) \cdot x_j &= \sum_{i=1}^n (\adj A)_{ji} \cdot b_i \\
  &= \sum_{i=1}^n (-1)^{i+j} b_i \det A(i|j) \\
  &= \det(A(j, b))
\end{align*}
where $A(j, b)$ is the matrix obtained from $A$ by replacing column $j$ with the vector $b$.

\begin{theorem}[Cramer's Rule]
  \label{thm:cramers_rule}
  Let $A \in \M_{n \times n}(K)$ with $\det A \neq 0$. Let $b = (b_1, \dots, b_n)\transp \in K^n$. Then, the system $A x = b$, where $x = (x_1, \dots, x_n)\transp$, has a unique solution, and it is given by the formula:
  \[
  x_j = \frac{\det A(j, b)}{\det A} \quad \text{for all } 1 \le j \le n,
  \]
  where $A(j,b)$ is the matrix obtained from $A$ by replacing column $j$ with $b$.
\end{theorem}

\begin{proof}[Proof of \Cref{thm:cramers_rule}]
  Since $\det A \neq 0$, the matrix $A$ is invertible by \cref{thm:inverse_adjoint}. Then $x := A^{-1} b$ solves the system, because $A(A^{-1}b) = b$, and it is the only solution: any $x$ with $Ax = b$ satisfies $x = A^{-1}(Ax) = A^{-1}b$. So, the system has a unique solution, and there is no need to appeal to Linear Algebra I for it. We have just shown above that $(\det A) \cdot x_j = \det A(j, b)$. Since $\det A \neq 0$, we can safely divide to obtain $x_j = \frac{\det A(j, b)}{\det A}$.
\end{proof}

\begin{example}[Solving a System using Cramer's Rule]
  Consider the system $\left(\begin{smallmatrix} 2 & 1 \\ 1 & 3 \end{smallmatrix}\right) \left(\begin{smallmatrix} x_1 \\ x_2 \end{smallmatrix}\right) = \left(\begin{smallmatrix} 5 \\ 5 \end{smallmatrix}\right)$. We have $\det(A) = 2(3) - 1(1) = 5$. Using Cramer's rule:
  \[
  x_1 = \frac{\det \begin{pmatrix} 5 & 1 \\ 5 & 3 \end{pmatrix}}{5} = \frac{15 - 5}{5} = 2, \quad x_2 = \frac{\det \begin{pmatrix} 2 & 5 \\ 1 & 5 \end{pmatrix}}{5} = \frac{10 - 5}{5} = 1.
  \]
  Thus, the solution is $x = \binom{2}{1}$.
\end{example}

\subsection{Determinants of Endomorphisms}

\begin{warmup} Let $A \in \M_{n \times n}(K)$, and let $P \in \GL(n, K)$ (meaning $P$ is an invertible $n \times n$ matrix). Then, by multiplicativity (\cref{thm:det_multiplicative}):
  \[
  \det(P^{-1} \cdot A \cdot P) = (\det P^{-1}) \cdot (\det A) \cdot (\det P) = \frac{1}{\det P} \cdot \det A \cdot \det P = \det A.
  \]
  In other words, similar matrices mathematically have the exact same determinant.
\end{warmup}

Let $V$ be a finite-di\-men\-sional vector space over $K$, and put $n := \dim V$. Let $T : V \to V$ be a linear map (an endomorphism). Pick a basis $\mathcal{B}$ for $V$. We have a matrix representation $[T]_{\mathcal{B}}^{\mathcal{B}} \in \M_{n \times n}(K)$, and we can consider its determinant $\det [T]_{\mathcal{B}}^{\mathcal{B}} \in K$.

\begin{question} Does $\det [T]_{\mathcal{B}}^{\mathcal{B}}$ depend on the specific choice of the basis $\mathcal{B}$?
\end{question}

\begin{answer} Let $\mathcal{B}'$ be another arbitrary basis for $V$. Denote by $P := [\id_V]_{\mathcal{B}}^{\mathcal{B}'}$ the transition matrix between the basis $\mathcal{B}'$ and $\mathcal{B}$. Recall that $[T]_{\mathcal{B}'}^{\mathcal{B}'} = [\id_V]_{\mathcal{B}'}^{\mathcal{B}} \cdot [T]_{\mathcal{B}}^{\mathcal{B}} \cdot [\id_V]_{\mathcal{B}}^{\mathcal{B}'}$.
  We also formally have $[\id_V]_{\mathcal{B}'}^{\mathcal{B}} = ([\id_V]_{\mathcal{B}}^{\mathcal{B}'})^{-1} = P^{-1}$.
  So, $[T]_{\mathcal{B}'}^{\mathcal{B}'} = P^{-1} \cdot [T]_{\mathcal{B}}^{\mathcal{B}} \cdot P$. This implies $\det [T]_{\mathcal{B}'}^{\mathcal{B}'} = \det [T]_{\mathcal{B}}^{\mathcal{B}}$.

  Recall also that if $T, S : V \to V$ are two linear maps, then $[S \circ T]_{\mathcal{B}}^{\mathcal{B}} = [S]_{\mathcal{B}}^{\mathcal{B}} \cdot [T]_{\mathcal{B}}^{\mathcal{B}}$. If $T$ is an isomorphism, then $[T^{-1}]_{\mathcal{B}}^{\mathcal{B}} = ([T]_{\mathcal{B}}^{\mathcal{B}})^{-1}$. Finally, $[\id_V]_{\mathcal{B}}^{\mathcal{B}} = I$.
\end{answer}

\begin{corollary}[Determinant of an Endomorphism]
  \label{cor:det_endomorphism}
  Let $V$ be a finite-di\-men\-sional vector space over $K$, and let $T : V \to V$ be a linear map. Then, $\det([T]_{\mathcal{B}}^{\mathcal{B}})$ does NOT depend on the choice of the basis $\mathcal{B}$ of $V$. We denote this intrinsic scalar simply by $\det(T) \in K$.
\end{corollary}

% Creator: Opus 5
% The notes list these four properties in running prose, with no environment and
% no proof, and ch. 21 uses the second of them to characterise eigenvalues.
% Unnumbered, so that 20.c.1--20.c.8 do not move; cite it with \cpageref.
\begin{proposition*}[Properties of the Determinant of an Endomorphism]
\label{prop:det_endomorphism_properties}
Let $V$ be a vector space over $K$ with $1 \le \dim V < \infty$. The function $\det : \End(V) \to K$, $T \mapsto \det(T)$, has the following elegant properties:
\begin{enumerate}[label=\textbf{(\alph*)}]
  \item $\det(S \circ T) = \det(S) \cdot \det(T)$ for all $S, T \in \End(V)$.
  \item $T$ is an isomorphism if and only if $\det T \neq 0$. If $T$ is an isomorphism, then $\det(T^{-1}) = \frac{1}{\det T}$.
  \item $\det(\id_V) = 1$.
  \item $\det(0) = 0 \in K$ (where $0$ is the zero linear map).
\end{enumerate}
\end{proposition*}

\begin{proof}
Fix a basis $\mathcal{B}$ of $V$ and write $n := \dim V$, so that $\det(T) = \det [T]_{\mathcal{B}}^{\mathcal{B}}$ by \cref{cor:det_endomorphism}. Each statement is then the corresponding statement about matrices.

\textbf{(a)} The composition rule gives $[S \circ T]_{\mathcal{B}}^{\mathcal{B}} = [S]_{\mathcal{B}}^{\mathcal{B}} \cdot [T]_{\mathcal{B}}^{\mathcal{B}}$, and \cref{thm:det_multiplicative} turns the product of matrices into the product of determinants.

\textbf{(b)} By \cref{prop:representation_of_isomorphisms}, $T$ is an isomorphism precisely when $[T]_{\mathcal{B}}^{\mathcal{B}}$ is invertible, which by \cref{thm:inverse_adjoint} happens precisely when $\det [T]_{\mathcal{B}}^{\mathcal{B}} \neq 0$. The formula for $\det(T^{-1})$ is then \cref{cor:det_inverse}.

\textbf{(c)} $[\id_V]_{\mathcal{B}}^{\mathcal{B}} = I_n$ and $\det(I_n) = 1$, the latter being the normalisation in \cref{def:determinant_func}.

\textbf{(d)} $[0]_{\mathcal{B}}^{\mathcal{B}}$ is the zero matrix, and every summand of the Leibniz formula of \cref{thm:unique_det} then carries a factor $0$. This is where $\dim V \ge 1$ is needed: for $n \ge 1$ each summand has at least one factor. \qedhere
\end{proof}

\begin{importantremark}
  \label{rem:det_needs_endomorphism}
  In contrast to an endomorphism $T : V \to V$, if we take a linear map $T : V \to W$ between two vector spaces of the same dimension and bases $\mathcal{B}, \mathcal{C}$ of $V, W$, then $\det[T]_{\mathcal{C}}^{\mathcal{B}}$ generally depends on the choice of $\mathcal{B}$ and $\mathcal{C}$.
\end{importantremark}

% Creator: Opus 5
\subsection{Solutions to the Exercises}
\label{sec:solutions_determinants}

\begin{exercisesolution}[Proof of \cref{exc:lin_dep_2x2}]
\label{sol:lin_dep_2x2}
Write $u := \binom{a}{c}$ and $v := \binom{b}{d}$ for the two columns. By \cref{prop:dependence_is_redundancy}, the list $(u,v)$ is linearly dependent exactly when one of the two vectors is a linear combination of the other. Explicitly: dependence gives scalars $\mu, \nu \in K$, not both zero, with $\mu u + \nu v = 0$. If $\mu \neq 0$ then $u = -\frac{\nu}{\mu} v$, which is the first alternative with $\lambda := -\frac{\nu}{\mu}$; and if $\mu = 0$ then $\nu \neq 0$, so $v = 0 = 0 \cdot u$, which is the second alternative with $\tau := 0$. Conversely, either alternative exhibits a non-trivial vanishing combination, so the columns are dependent. Note that the case $u = 0$ is contained in the first alternative, taking $\lambda := 0$.
\end{exercisesolution}

\begin{exercisesolution}[Proof of \cref{exc:size_of_Sn}]
\label{sol:size_of_Sn}
A permutation $\sigma \in S_n$ is determined by the list $(\sigma(1), \dots, \sigma(n))$, which is an arrangement of $1, \dots, n$ with no repetitions. Building such a list, there are $n$ choices for $\sigma(1)$; once it is fixed, $\sigma(2)$ may be any of the remaining $n-1$ values; then $n-2$ for $\sigma(3)$, and so on, with a single choice left for $\sigma(n)$. Distinct sequences of choices give distinct permutations and every permutation arises exactly once, so
\[
|S_n| = n \cdot (n-1) \cdot (n-2) \cdots 2 \cdot 1 = n!. \qedhere
\]
\end{exercisesolution}

\begin{exercisesolution}[Proof of \cref{exc:block_lower}]
\label{sol:block_lower}
Take the transpose. By \cref{thm:det_transpose}, $\det \left(\begin{smallmatrix} A & 0 \\ B & C \end{smallmatrix}\right) = \det \left( \left(\begin{smallmatrix} A & 0 \\ B & C \end{smallmatrix}\right)\transp \right)$, and transposing a block matrix transposes each block and reflects the arrangement, giving
\[
\left(\begin{smallmatrix} A & 0 \\ B & C \end{smallmatrix}\right)\transp = \left(\begin{smallmatrix} A\transp & B\transp \\ 0 & C\transp \end{smallmatrix}\right),
\]
where now $A\transp \in \M_{r \times r}(K)$, $C\transp \in \M_{s \times s}(K)$ and $B\transp \in \M_{r \times s}(K)$. This is exactly the shape of \cref{prop:block_matrix}, so
\[
\det \begin{pmatrix} A & 0 \\ B & C \end{pmatrix} = \det \begin{pmatrix} A\transp & B\transp \\ 0 & C\transp \end{pmatrix} = \det(A\transp) \cdot \det(C\transp) = \det(A) \cdot \det(C),
\]
using \cref{thm:det_transpose} once more on each factor.
\end{exercisesolution}

% Official Solution (Problem Sheet 16, Exercise 2)
\begin{exercisesolution}[Solution to \cref{exc:invertibility_nilpotency_2x2}]
\label{sol:invertibility_nilpotency_2x2}
\begin{enumerate}[label=\textbf{(\alph*)}]
    \item \textbf{Invertibility condition:}
    Let $A = \begin{pmatrix} a & b \\ c & d \end{pmatrix}$.
    \begin{itemize}
        \item If $ad - bc \ne 0$, the matrix $B = \frac{1}{ad-bc} \begin{pmatrix} d & -b \\ -c & a \end{pmatrix}$ satisfies $AB = BA = I_2$, so $A$ is invertible.
        \item Conversely, if $A$ is invertible, there exists $A^{-1} \in \M_{2 \times 2}(K)$ such that $A A^{-1} = I_2$. Taking determinants gives $\det(A) \det(A^{-1}) = \det(I_2) = 1$, which forces $\det(A) = ad - bc \ne 0$.
    \end{itemize}

    \item \textbf{Nilpotency condition:}
    We compute the characteristic polynomial of $A$ in $K[t]$:
    \begin{align*}
    \det(tI_2 - A) &= \det \begin{pmatrix} t - a & -b \\ -c & t - d \end{pmatrix} = (t - a)(t - d) - bc \\
    &= t^2 - (a + d)t + (ad - bc) = t^2 - (\operatorname{tr} A)t + \det A.
    \end{align*}
    Thus $\det(tI_2 - A) = t^2 \iff \operatorname{tr} A = 0$ and $\det A = 0$.
    \begin{itemize}
        \item If $\operatorname{tr} A = 0$ and $\det A = 0$, direct computation yields:
        \[
        A^2 = \begin{pmatrix} a & b \\ c & d \end{pmatrix} \begin{pmatrix} a & b \\ c & -a \end{pmatrix} = \begin{pmatrix} a^2 + bc & b(a+d) \\ c(a+d) & bc + d^2 \end{pmatrix} = \begin{pmatrix} a^2 + bc & 0 \\ 0 & bc + a^2 \end{pmatrix} = \begin{pmatrix} 0 & 0 \\ 0 & 0 \end{pmatrix},
        \]
        since $a^2 + bc = -ad + bc = -(ad - bc) = -\det A = 0$.
        \item Conversely, if $A^2 = 0$, then $(\det A)^2 = \det(A^2) = \det(0) = 0 \implies \det A = 0$.
        By the Cayley-Hamilton relation for $2 \times 2$ matrices, $A^2 - (\operatorname{tr} A)A + (\det A)I_2 = 0$, which reduces to $(\operatorname{tr} A)A = 0$.
        If $A \ne 0$, this implies $\operatorname{tr} A = 0$. If $A = 0$, then $\operatorname{tr} A = 0$ and $\det A = 0$ trivially.
    \end{itemize}
    Therefore, $A^2 = 0 \iff \det(tI_2 - A) = t^2$. \qedhere
\end{enumerate}
\exinfo{Official Solution (Problem Sheet 16, Exercise 2). Proves the $2 \times 2$ determinant invertibility and nilpotency conditions.}
\end{exercisesolution}

% Official Solution (Problem Sheet 16, Exercise 4)
\begin{exercisesolution}[Solution to \cref{exc:monomial_matrix_multilinear}]
\label{sol:monomial_matrix_multilinear}
Let $D(A) = A(j_1, k_1) A(j_2, k_2) \cdots A(j_n, k_n)$ where $j_i, k_i \in \{1, \dots, n\}$.

\begin{itemize}
    \item \textbf{($\Leftarrow$) If $j_1, \dots, j_n$ are pairwise distinct:}
    Then $\{j_1, \dots, j_n\} = \{1, \dots, n\}$. For each row index $r \in \{1, \dots, n\}$, exactly one factor $A(r, k_{\ell})$ with $j_{\ell} = r$ appears in the product.
    Holding all rows fixed except row $r$, $D(A) = C \cdot A(r, k_{\ell})$ where $C = \prod_{m \ne \ell} A(j_m, k_m)$ is a constant independent of row $r$.
    The map $A \mapsto C \cdot A(r, k_{\ell})$ is linear in row $r$, so $D$ is $n$-linear.

    \item \textbf{($\Rightarrow$) If $D$ is $n$-linear:}
    Suppose the indices $j_1, \dots, j_n$ are not pairwise distinct.
    Then by the Pigeonhole Principle, there exists at least one row index $r \in \{1, \dots, n\}$ that appears $m \ge 2$ times in the list $(j_1, \dots, j_n)$.
    Scaling row $r$ by a scalar $c \in K$ replaces each entry $A(r, \cdot)$ by $c A(r, \cdot)$, so:
    \[
    D(c R_r) = c^m D(A).
    \]
    On the other hand, $n$-linearity requires $D(c R_r) = c D(A)$ for all $c \in K$.
    Choosing $A$ to be the all-ones matrix $J_n$, we have $D(J_n) = 1 \ne 0$, which would require:
    \[
    c^m = c \qquad \text{for all } c \in K.
    \]
    Since $K \subseteq \mathbb{C}$ is a subfield of $\mathbb{C}$, it is an infinite field containing $\mathbb{Q}$.
    The polynomial $t^m - t = 0$ has at most $m$ roots in $K$, so $c^m = c$ cannot hold for all $c \in K$ (for example, taking $c = 2$ gives $2^m \ne 2$ since $m \ge 2$).
    This contradiction shows that the row indices $j_1, \dots, j_n$ must be pairwise distinct. \qedhere
\end{itemize}
\exinfo{Official Solution (Problem Sheet 16, Exercise 4). Characterizes multilinearity of monomial expressions in matrix entries.}
\end{exercisesolution}

% Official Solution (Problem Sheet 16, Exercise 3)
\begin{exercisesolution}[Solution to \cref{exc:multilinear_candidates_3x3}]
\label{sol:multilinear_candidates_3x3}
\begin{enumerate}[label=\textbf{(\alph*)}]
    \item $D(A) = A_{11} + A_{22} + A_{33}$ (\textbf{Not $3$-linear}):
    For $A = I_3$, $D(I_3) = 3$. Scaling the first row by $c = 2$ gives $D(2R_1, R_2, R_3) = 2(1) + 1 + 1 = 4 \ne 2 \cdot D(I_3) = 6$.

    \item $D(A) = A_{11}^2 + 3A_{11} A_{22}$ (\textbf{Not $3$-linear}):
    The term $A_{11}^2$ is quadratic in row $1$, so scaling row $1$ by $c$ scales $A_{11}^2$ by $c^2 \ne c$. Moreover, row $3$ does not appear.

    \item $D(A) = A_{11} A_{12} A_{33}$ (\textbf{Not $3$-linear}):
    Row $1$ appears twice (degree $2$ in row $1$), while row $2$ is completely missing.

    \item $D(A) = A_{13} A_{22} A_{32} + 5A_{12} A_{22} A_{32} = (A_{13} + 5A_{12}) A_{22} A_{32}$ (\textbf{Is $3$-linear}):
    It is the product of three linear forms in row $1$, row $2$, and row $3$, respectively.

    \item $D(A) = 0$ (\textbf{Is $3$-linear}):
    The zero map trivially satisfies all $3$-linearity conditions.

    \item $D(A) = 1$ (\textbf{Not $3$-linear}):
    Scaling row $1$ by $c = 0$ gives $D(0, R_2, R_3) = 1 \ne 0 \cdot D(A) = 0$. \qedhere
\end{enumerate}
\exinfo{Official Solution (Problem Sheet 16, Exercise 3). Identifies which functions on $3 \times 3$ matrices are 3-linear.}
\end{exercisesolution}

% Official Solution (Problem Sheet 16, Exercise 1)
\begin{exercisesolution}[Solution to \cref{exc:det_2x2_alternating_bilinear}]
\label{sol:det_2x2_alternating_bilinear}
Let $A = \begin{pmatrix} a & b \\ c & d \end{pmatrix} \in \M_{2 \times 2}(K)$, so $\det A = ad - bc$.

\begin{itemize}
    \item \textbf{Linearity in rows:}
    Fixing $(c, d)$, the map $(a, b) \mapsto ad - bc = d \cdot a - c \cdot b$ is linear in the first row.
    Fixing $(a, b)$, the map $(c, d) \mapsto ad - bc = -b \cdot c + a \cdot d$ is linear in the second row.
    If the two rows are equal, then $(a, b) = (c, d)$, so $\det A = ac - bc = ac - ac = 0$, so $\det$ is alternating in rows.

    \item \textbf{Linearity in columns:}
    Fixing column $2$, $(a, c) \mapsto ad - bc = d \cdot a - b \cdot c$ is linear in column $1$.
    Fixing column $1$, $(b, d) \mapsto ad - bc = -c \cdot b + a \cdot d$ is linear in column $2$.
    If the two columns are equal, then $(a, c) = (b, d)$, so $\det A = ab - ba = 0$, so $\det$ is alternating in columns. \qedhere
\end{itemize}
\exinfo{Official Solution (Problem Sheet 16, Exercise 1). Proves the $2 \times 2$ determinant is an alternating bilinear function in rows and columns.}
\end{exercisesolution}

% Official Solution (Problem Sheet 16, Exercise 5)
\begin{exercisesolution}[Solution to \cref{exc:laplace_expansion_3x3_alternating_trilinear}]
\label{sol:laplace_expansion_3x3_alternating_trilinear}
Let $A = (C_1, C_2, C_3) \in \M_{3 \times 3}(K)$. Expanding the $2 \times 2$ determinants in the formula:
\[
D(A) = A_{11}(A_{22}A_{33} - A_{23}A_{32}) - A_{12}(A_{21}A_{33} - A_{23}A_{31}) + A_{13}(A_{21}A_{32} - A_{22}A_{31}).
\]
\begin{itemize}
    \item \textbf{$3$-linearity in columns:}
    In every term of the expansion, each column index $1, 2, 3$ appears exactly once as the second index of an entry.
    Hence $D(A)$ is linear in $C_1$ (holding $C_2, C_3$ fixed), linear in $C_2$, and linear in $C_3$.
    \item \textbf{Alternating property:}
    \begin{itemize}
        \item If $C_1 = C_2$, then $A_{i1} = A_{i2}$ for all $i \in \{1, 2, 3\}$.
        The third term vanishes because $A_{21}A_{32} - A_{22}A_{31} = A_{21}A_{31} - A_{21}A_{31} = 0$.
        The first two terms become identical with opposite signs:
        $A_{11}(A_{21}A_{33} - A_{23}A_{31}) - A_{11}(A_{21}A_{33} - A_{23}A_{31}) = 0$.
        \item If $C_1 = C_3$ or $C_2 = C_3$, a completely symmetric cancellation occurs.
    \end{itemize}
\end{itemize}
Thus $D$ is an alternating $3$-linear function of the columns.
\exinfo{Official Solution (Problem Sheet 16, Exercise 5). Demonstrates the alternating trilinear property of the $3 \times 3$ Laplace expansion.}
\end{exercisesolution}

% Solution: Gemini / Official Hybrid (Problem Sheet 16, Lecture)
\begin{exercisesolution}[Solution to \cref{exc:triangular_matrix_determinant}]
\label{sol:triangular_matrix_determinant}
Let $A \in \M_{n \times n}(K)$ be an upper triangular matrix ($A_{ij} = 0$ for all $i > j$).
By the Leibniz formula (\cref{thm:unique_det}):
\[
\det(A) = \sum_{\sigma \in S_n} \sgn(\sigma) \prod_{i=1}^n A_{i, \sigma(i)}.
\]
Suppose $\sigma \in S_n$ is a permutation with $\sigma \ne \id$.
Let $k \in \{1, \dots, n\}$ be the smallest index such that $\sigma(k) \ne k$.
Since $\sigma(1) = 1, \dots, \sigma(k-1) = k-1$, we must have $\sigma(k) > k$.
Because $\sigma \colon \{1, \dots, n\} \to \{1, \dots, n\}$ is a bijection, the remaining elements $\{\sigma(k+1), \dots, \sigma(n)\}$ must be mapped into $\{1, \dots, n\} \setminus \{1, \dots, k-1, \sigma(k)\}$.
Since there are $n - k$ remaining positions and only $n - k - 1$ values greater than $k$, by the Pigeonhole Principle there must exist some index $j > k$ such that $\sigma(j) \le k < j$.
For this index $j$, we have $j > \sigma(j)$, which implies $A_{j, \sigma(j)} = 0$.
Consequently, the product $\prod_{i=1}^n A_{i, \sigma(i)} = 0$ for every $\sigma \ne \id$.
The only non-zero term is $\sigma = \id$, which yields:
\[
\det(A) = \sgn(\id) \prod_{i=1}^n A_{ii} = A_{11} A_{22} \cdots A_{nn}.
\]
For a lower triangular matrix, transposing gives an upper triangular matrix with the same diagonal, so $\det(A) = \det(A\transp) = A_{11} \cdots A_{nn}$.
\exinfo{Hybrid Solution (Gemini 3.7 Flash / Official Problem Sheet 16). Proves the determinant of a triangular matrix is the product of its diagonal entries.}
\end{exercisesolution}

% Solution: Gemini / Official Hybrid (Problem Sheet 16, Combinatorics)
\begin{exercisesolution}[Solution to \cref{exc:symmetric_group_s4_and_alternating}]
\label{sol:symmetric_group_s4_and_alternating}
\begin{enumerate}[label=\textbf{(\alph*)}]
    \item The group $S_4$ contains $4! = 24$ permutations, classified by cycle structure:
    \begin{itemize}
        \item \textbf{Identity ($1$ element, $\sgn = +1$):} $\id$.
        \item \textbf{$2$-cycles / Transpositions ($6$ elements, $\sgn = -1$):}
        $(1\,2), (1\,3), (1\,4), (2\,3), (2\,4), (3\,4)$.
        \item \textbf{$3$-cycles ($8$ elements, $\sgn = +1$):} \\
        $(1\,2\,3), (1\,3\,2), (1\,2\,4), (1\,4\,2), (1\,3\,4), (1\,4\,3), (2\,3\,4), (2\,4\,3)$.
        \item \textbf{$4$-cycles ($6$ elements, $\sgn = -1$):}
        $(1\,2\,3\,4), (1\,2\,4\,3), (1\,3\,2\,4), (1\,3\,4\,2), (1\,4\,2\,3), (1\,4\,3\,2)$.
        \item \textbf{Double transpositions ($3$ elements, $\sgn = +1$):}
        $(1\,2)(3\,4), (1\,3)(2\,4), (1\,4)(2\,3)$.
    \end{itemize}
    There are $1 + 8 + 3 = 12$ even permutations and $6 + 6 = 12$ odd permutations.
    The complete Leibniz formula for $A \in \M_{4 \times 4}(K)$ is:
    \begin{align*}
    \det(A) &= A_{11}A_{22}A_{33}A_{44} + A_{12}A_{23}A_{31}A_{44} + A_{13}A_{21}A_{32}A_{44} + A_{12}A_{24}A_{33}A_{41} \\
    &\quad + A_{14}A_{21}A_{33}A_{42} + A_{13}A_{24}A_{31}A_{42} + A_{14}A_{22}A_{31}A_{43} + A_{11}A_{23}A_{34}A_{42} \\
    &\quad + A_{11}A_{24}A_{32}A_{43} + A_{12}A_{21}A_{34}A_{43} + A_{13}A_{22}A_{34}A_{41} + A_{14}A_{23}A_{32}A_{41} \\
    &\quad - A_{12}A_{21}A_{33}A_{44} - A_{13}A_{22}A_{31}A_{44} - A_{14}A_{22}A_{33}A_{41} - A_{11}A_{23}A_{32}A_{44} \\
    &\quad - A_{11}A_{24}A_{33}A_{42} - A_{11}A_{22}A_{34}A_{43} - A_{12}A_{23}A_{34}A_{41} - A_{12}A_{24}A_{31}A_{43} \\
    &\quad - A_{13}A_{21}A_{34}A_{42} - A_{13}A_{24}A_{32}A_{41} - A_{14}A_{21}A_{32}A_{43} - A_{14}A_{23}A_{31}A_{42}.
    \end{align*}

    \item For $n = 1$: $|S_1| = 1$, so there is $1$ even permutation ($\id$).
    For $n \ge 2$: Fix the transposition $\tau := (1\,2) \in S_n$.
    The map $\Phi \colon A_n \to S_n \setminus A_n$ given by $\Phi(\sigma) := \tau \circ \sigma$ is a bijection, since $\sgn(\tau \circ \sigma) = \sgn(\tau) \sgn(\sigma) = -\sgn(\sigma)$, and $\Phi$ is its own inverse.
    Thus the number of even permutations in $S_n$ is exactly:
    \[
    |A_n| = \frac{|S_n|}{2} = \frac{n!}{2}. \qedhere
    \]
\end{enumerate}
\exinfo{Hybrid Solution (Gemini 3.7 Flash / Official Problem Sheet 16). Classifies permutations in $S_4$, writes the 24-term Leibniz formula, and proves $|A_n| = n!/2$.}
\end{exercisesolution}

% Solution: Gemini / Official Hybrid (Problem Sheet 16 / Theory)
\begin{exercisesolution}[Solution to \cref{exc:multiplicativity_via_leibniz}]
\label{sol:multiplicativity_via_leibniz}
Let $A, B \in \M_{n \times n}(K)$. Write $B$ in terms of its columns $B = (b_1, \dots, b_n)$ where $b_j = \sum_{s_j=1}^n B(s_j, j) e_{s_j}$.
Then the columns of $AB$ are $A b_j = \sum_{s_j=1}^n B(s_j, j) A e_{s_j}$.
Using the multilinearity of $\det$ with respect to columns:
\[
\det(AB) = \det(A b_1, \dots, A b_n) = \sum_{s_1=1}^n \cdots \sum_{s_n=1}^n \left(\prod_{j=1}^n B(s_j, j)\right) \det(A e_{s_1}, \dots, A e_{s_n}).
\]
Because $\det$ is alternating in columns, $\det(A e_{s_1}, \dots, A e_{s_n}) = 0$ whenever two of the indices $s_1, \dots, s_n$ coincide.
Thus non-zero summands correspond precisely to permutations $\sigma \in S_n$ with $(s_1, \dots, s_n) = (\sigma(1), \dots, \sigma(n))$.
For any $\sigma \in S_n$, reordering the columns $(A e_{\sigma(1)}, \dots, A e_{\sigma(n)})$ back to $(A e_1, \dots, A e_n)$ requires transpositions that produce a factor of $\sgn(\sigma)$:
\[
\det(A e_{\sigma(1)}, \dots, A e_{\sigma(n)}) = \sgn(\sigma) \det(A e_1, \dots, A e_n) = \sgn(\sigma) \det(A).
\]
Substituting this back gives:
\[
\det(AB) = \det(A) \sum_{\sigma \in S_n} \sgn(\sigma) \prod_{j=1}^n B(\sigma(j), j) = \det(A) \det(B\transp) = \det(A) \det(B). \qedhere
\]
\exinfo{Hybrid Solution (Gemini 3.7 Flash / Biran Lecture Notes). Derives determinant multiplicativity directly from column Leibniz expansion.}
\end{exercisesolution}

% Official Solution (Problem Sheet 16, Multiple Choice)
\begin{exercisesolution}[Solution to \cref{exc:general_determinant_properties_mc}]
\label{sol:general_determinant_properties_mc}
\begin{itemize}
    \item \textbf{(a) True:} Multiplicativity of the determinant (\cref{thm:det_multiplicative}).
    \item \textbf{(b) True:} $\det A \ne 0 \iff A$ is invertible $\iff$ columns are linearly independent.
    \item \textbf{(c) True:} $\det(AB) = \det(A)\det(B) = \det(B)\det(A) = \det(BA)$ in the commutative field $\mathbb{R}$.
    \item \textbf{(d) False:} For $n \ge 2$, $n$-linearity implies $\det(\lambda A) = \lambda^n \det(A) \ne \lambda \det(A)$ in general.
    \item \textbf{(e) False:} The determinant is multilinear, not linear; for example $\det(I_2 + I_2) = \det(2I_2) = 4 \ne 1 + 1 = 2$.
\end{itemize}
The correct statements are \textbf{(a), (b), and (c)}.
\exinfo{Official Solution (Problem Sheet 16, Multiple Choice). Evaluates general determinant properties including multiplicativity and scaling.}
\end{exercisesolution}

% Official Solution (Problem Sheet 16, Single Choice)
\begin{exercisesolution}[Solution to \cref{exc:determinant_properties_linearity_sc}]
\label{sol:determinant_properties_linearity_sc}
Statement \textbf{(c)} is false: the determinant function $\det \colon \M_{n \times n}(K) \to K$ is \emph{not} a linear map for $n \ge 2$, since $\det(A + B) \ne \det(A) + \det(B)$ in general (e.g. $\det(I_n + I_n) = 2^n \ne 1 + 1 = 2$).
The remaining statements are all true:
\begin{itemize}
    \item (a) is the fundamental invertibility theorem (\cref{thm:inverse_adjoint}).
    \item (b) is the triangular matrix formula (\cref{exc:triangular_matrix_determinant}).
    \item (d) is true since $\det(\operatorname{diag}(\lambda, 1, \dots, 1)) = \lambda$ for any $\lambda \in K$.
\end{itemize}
Therefore, the false statement is \textbf{(c)}.
\exinfo{Official Solution (Problem Sheet 16, Single Choice). Tests understanding of non-linearity and fundamental properties of the determinant.}
\end{exercisesolution}

% Official Solution (Problem Sheet 17, Exercise 2)
\begin{exercisesolution}[Solution to \cref{exc:row_operations_and_determinant_divisibility}]
\label{sol:row_operations_and_determinant_divisibility}
\begin{enumerate}[label=\textbf{(\alph*)}]
    \item \leavevmode
    \begin{enumerate}[label=\textbf{(\roman*)}]
        \item By $n$-linearity on the $i$-th row, $\det(R_1, \dots, \lambda R_i, \dots, R_n) = \lambda \det(R_1, \dots, R_i, \dots, R_n) = \lambda \det A$.
        \item By the alternating property (\cref{prop:alternating_swap} \textbf{(b)}), swapping any two rows flips the sign of the determinant: $\det B = -\det A$.
        \item Expanding by linearity on row $j$: $\det(\dots, R_j + \lambda R_i, \dots) = \det(\dots, R_j, \dots) + \lambda \det(\dots, R_i, \dots) = \det A + 0 = \det A$, since the second matrix has two identical rows $R_i$.
    \end{enumerate}

    \item Let $M = \begin{pmatrix} 2 & 1 & 3 & 6 \\ 0 & 4 & 7 & 9 \\ 1 & 8 & 1 & 9 \\ 4 & 4 & 0 & 6 \end{pmatrix}$.
    Perform the Type I row operation $R_4 \to 1000 R_1 + 100 R_2 + 10 R_3 + R_4$. By part \textbf{(a)(iii)}, this preserves the determinant:
    \[
    \det(M) = \det \begin{pmatrix}
      2 & 1 & 3 & 6 \\
      0 & 4 & 7 & 9 \\
      1 & 8 & 1 & 9 \\
      2014 & 1484 & 3710 & 6996
    \end{pmatrix}.
    \]
    Since $2014 = 106 \times 19$, $1484 = 106 \times 14$, $3710 = 106 \times 35$, and $6996 = 106 \times 66$, we factor out $106$ from the fourth row using part \textbf{(a)(i)}:
    \[
    \det(M) = 106 \cdot \det \begin{pmatrix}
      2 & 1 & 3 & 6 \\
      0 & 4 & 7 & 9 \\
      1 & 8 & 1 & 9 \\
      19 & 14 & 35 & 66
    \end{pmatrix} \in 106 \mathbb{Z}.
    \]
    Thus $\det(M)$ is divisible by $106$. \qedhere
\end{enumerate}
\exinfo{Official Solution (Problem Sheet 17, Exercise 2). Connects elementary row operations to determinant scaling and proves 106-divisibility without calculation.}
\end{exercisesolution}

% Official Solution (Problem Sheet 17, Single Choice)
\begin{exercisesolution}[Solution to \cref{exc:invariance_under_matrix_operations_sc}]
\label{sol:invariance_under_matrix_operations_sc}
Interchanging two rows flips the sign of the determinant ($\det \to -\det$), so the determinant in general changes.
The other operations strictly preserve the determinant:
\begin{itemize}
    \item (b) Adding a scalar multiple of one row to another preserves $\det$ (\cref{thm:row_operations} \textbf{(a)}).
    \item (c) $\det(A\transp) = \det(A)$ (\cref{thm:det_transpose}).
    \item (d) $\det(P^{-1}AP) = \det(A)$ (\cref{cor:det_endomorphism}).
\end{itemize}
Therefore, the correct answer is \textbf{(a)}.
\exinfo{Official Solution (Problem Sheet 17, Single Choice). Tests invariance of determinants under transposition, similarity, and row operations.}
\end{exercisesolution}

% Official Solution (Problem Sheet 17, Multiple Choice)
\begin{exercisesolution}[Solution to \cref{exc:row_transformations_2x2_mc}]
\label{sol:row_transformations_2x2_mc}
Given $\det \begin{pmatrix} a & b \\ c & d \end{pmatrix} = 4$:
\begin{itemize}
    \item \textbf{(a) Incorrect:} $\det \begin{pmatrix} 2a & 2b \\ 2c & 2d \end{pmatrix} = 2^2 \det \begin{pmatrix} a & b \\ c & d \end{pmatrix} = 4 \times 4 = 16 \ne 8$.
    \item \textbf{(b) Correct:} Applying $R_2 \to R_2 - R_1$ preserves the determinant, so the value remains $4$.
    \item \textbf{(c) Correct:} Applying $R_2 \to R_2 + 2R_1$ preserves the determinant, so the value remains $4$.
    \item \textbf{(d) Correct:} Scaling the second row by $3$ scales the determinant by $3$, giving $3 \times 4 = 12$.
\end{itemize}
The correct statements are \textbf{(b), (c), and (d)}.
\exinfo{Official Solution (Problem Sheet 17, Multiple Choice). Evaluates determinant effects of row operations on a $2 \times 2$ matrix.}
\end{exercisesolution}

% Official Solution (Problem Sheet 17, Exercise 3)
\begin{exercisesolution}[Solution to \cref{exc:block_matrix_commuting_blocks}]
\label{sol:block_matrix_commuting_blocks}
Let $M = \begin{pmatrix} A & B \\ C & D \end{pmatrix}$ with $AC = CA$ and $\det A \ne 0$.
Consider the block matrix $N = \begin{pmatrix} I_n & 0 \\ -C & A \end{pmatrix}$.
Multiplying $N$ and $M$ blockwise:
\[
NM = \begin{pmatrix} I_n & 0 \\ -C & A \end{pmatrix} \begin{pmatrix} A & B \\ C & D \end{pmatrix} = \begin{pmatrix} A & B \\ -CA + AC & -CB + AD \end{pmatrix} = \begin{pmatrix} A & B \\ 0 & AD - CB \end{pmatrix},
\]
where the lower left block vanishes because $AC = CA$.
Taking determinants on both sides and using multiplicativity (\cref{thm:det_multiplicative}) and the block determinant formula (\cref{prop:block_matrix} and \cref{exc:block_lower}):
\[
\det(NM) = \det(N) \det(M) = \det(I_n) \det(A) \det(M) = \det(A) \det(M).
\]
On the other hand:
\[
\det(NM) = \det \begin{pmatrix} A & B \\ 0 & AD - CB \end{pmatrix} = \det(A) \det(AD - CB).
\]
Since $\det A \ne 0$, dividing both sides by $\det(A)$ yields:
\[
\det \begin{pmatrix} A & B \\ C & D \end{pmatrix} = \det(AD - CB). \qedhere
\]
\exinfo{Official Solution (Problem Sheet 17, Exercise 3). Proves block determinant identity $\det(M) = \det(AD - CB)$ when $AC = CA$ and $A$ is invertible.}
\end{exercisesolution}

% Official Solution (Problem Sheet 17, Exercise 4)
\begin{exercisesolution}[Solution to \cref{exc:determinants_via_gaussian_elimination}]
\label{sol:determinants_via_gaussian_elimination}
\begin{enumerate}[label=\textbf{(\alph*)}]
    \item \textbf{Matrix $A$:}
    Interchange $R_1 \leftrightarrow R_3$ (one swap, sign $-1$):
    \[
    \begin{pmatrix}
      1 & 2 & -3 & 1 \\
      -1 & 0 & 1 & 0 \\
      0 & 1 & 0 & 0 \\
      0 & -4 & 2 & -1
    \end{pmatrix}
    \xrightarrow{R_2 + R_1 \to R_2}
    \begin{pmatrix}
      1 & 2 & -3 & 1 \\
      0 & 2 & -2 & 1 \\
      0 & 1 & 0 & 0 \\
      0 & -4 & 2 & -1
    \end{pmatrix}
    \xrightarrow{R_4 + 2R_2 \to R_4}
    \begin{pmatrix}
      1 & 2 & -3 & 1 \\
      0 & 2 & -2 & 1 \\
      0 & 1 & 0 & 0 \\
      0 & 0 & -2 & 1
    \end{pmatrix}.
    \]
    Eliminating row $3$ using $R_3 \to 2R_3 - R_2$ gives $(0, 0, 2, -1)$, which is the exact negation of row $4$. Adding row $4$ yields a zero row, so $\rank(A) < 4$ and:
    \[
    \det(A) = 0.
    \]

    \item \textbf{Matrix $B$:}
    Expanding along the first row (where only $B_{15} = 1$ is non-zero):
    \[
    \det(B) = 1 \cdot (-1)^{1+5} \det \begin{pmatrix}
      -1 & 2 & 0 & 0 \\
      0 & 1 & 0 & -1 \\
      0 & 0 & -1 & 0 \\
      0 & 0 & 0 & 1
    \end{pmatrix}.
    \]
    This $4 \times 4$ submatrix is upper triangular, so its determinant is the product of diagonal entries $(-1)(1)(-1)(1) = 1$.
    Thus:
    \[
    \det(B) = 1.
    \]

    \item \textbf{Matrix $C$:}
    Perform Gaussian elimination to bring $C$ to upper triangular form:
    \[
    \begin{pmatrix}
      2 & -3 & 5 & 1 & 4 \\
      2 & -3 & 1 & -6 & 18 \\
      4 & -3 & 9 & 6 & 10 \\
      -2 & 4 & -6 & -1 & -1 \\
      -6 & 11 & -23 & -14 & 9
    \end{pmatrix}
    \xrightarrow{\substack{R_2 - R_1 \to R_2 \\ R_3 - 2R_1 \to R_3 \\ R_4 + R_1 \to R_4 \\ R_5 + 3R_1 \to R_5}}
    \begin{pmatrix}
      2 & -3 & 5 & 1 & 4 \\
      0 & 0 & -4 & -7 & 14 \\
      0 & 3 & -1 & 4 & 2 \\
      0 & 1 & -1 & 0 & 3 \\
      0 & 2 & -8 & -11 & 21
    \end{pmatrix}.
    \]
    Swap $R_2 \leftrightarrow R_4$ (one swap, sign $-1$):
    \[
    \xrightarrow{R_2 \leftrightarrow R_4}
    \begin{pmatrix}
      2 & -3 & 5 & 1 & 4 \\
      0 & 1 & -1 & 0 & 3 \\
      0 & 3 & -1 & 4 & 2 \\
      0 & 0 & -4 & -7 & 14 \\
      0 & 2 & -8 & -11 & 21
    \end{pmatrix}
    \xrightarrow{\substack{R_3 - 3R_2 \to R_3 \\ R_5 - 2R_2 \to R_5}}
    \begin{pmatrix}
      2 & -3 & 5 & 1 & 4 \\
      0 & 1 & -1 & 0 & 3 \\
      0 & 0 & 2 & 4 & -7 \\
      0 & 0 & -4 & -7 & 14 \\
      0 & 0 & -6 & -11 & 15
    \end{pmatrix}.
    \]
    Eliminate the third column below the diagonal using $R_3$:
    \[
    \xrightarrow{\substack{R_4 + 2R_3 \to R_4 \\ R_5 + 3R_3 \to R_5}}
    \begin{pmatrix}
      2 & -3 & 5 & 1 & 4 \\
      0 & 1 & -1 & 0 & 3 \\
      0 & 0 & 2 & 4 & -7 \\
      0 & 0 & 0 & 1 & 0 \\
      0 & 0 & 0 & 1 & -6
    \end{pmatrix}
    \xrightarrow{R_5 - R_4 \to R_5}
    \begin{pmatrix}
      2 & -3 & 5 & 1 & 4 \\
      0 & 1 & -1 & 0 & 3 \\
      0 & 0 & 2 & 4 & -7 \\
      0 & 0 & 0 & 1 & 0 \\
      0 & 0 & 0 & 0 & -6
    \end{pmatrix}.
    \]
    The matrix is now upper triangular with diagonal $(2, 1, 2, 1, -6)$.
    Accounting for the single row interchange:
    \[
    \det(C) = (-1) \cdot \big(2 \cdot 1 \cdot 2 \cdot 1 \cdot (-6)\big) = (-1) \cdot (-24) = 24. \qedhere
    \]
\end{enumerate}
\exinfo{Official Solution (Problem Sheet 17, Exercise 4). Computes determinants of three $4 \times 4$ and $5 \times 5$ matrices using Gaussian elimination.}
\end{exercisesolution}

% Official Solution (Problem Sheet 17, Exercise 5)
\begin{exercisesolution}[Solution to \cref{exc:vandermonde_determinant}]
\label{sol:vandermonde_determinant}
Let $V_n(x_1, \dots, x_n) = \det(x_i^{j-1})_{1 \le i, j \le n}$. We proceed by induction on $n \ge 2$.
\begin{itemize}
    \item \textbf{Base case $n = 2$:}
    $\det \begin{pmatrix} 1 & x_1 \\ 1 & x_2 \end{pmatrix} = x_2 - x_1$, which matches $\prod_{1 \le i < j \le 2} (x_j - x_i)$.
    \item \textbf{Induction step:}
    Assume the formula holds for $n-1$. For the $n \times n$ matrix, perform column operations from right to left:
    \[
    C_j \longrightarrow C_j - x_1 C_{j-1} \qquad \text{for } j = n, n-1, \dots, 2.
    \]
    This clears all entries of the first row except the top-left $1$:
    \[
    \begin{pmatrix}
      1 & 0 & 0 & \cdots & 0 \\
      1 & x_2 - x_1 & x_2(x_2 - x_1) & \cdots & x_2^{n-2}(x_2 - x_1) \\
      \vdots & \vdots & \vdots & \ddots & \vdots \\
      1 & x_n - x_1 & x_n(x_n - x_1) & \cdots & x_n^{n-2}(x_n - x_1)
    \end{pmatrix}.
    \]
    Expanding along the first row:
    \[
    V_n(x_1, \dots, x_n) = \det \begin{pmatrix}
      x_2 - x_1 & x_2(x_2 - x_1) & \cdots & x_2^{n-2}(x_2 - x_1) \\
      \vdots & \vdots & \ddots & \vdots \\
      x_n - x_1 & x_n(x_n - x_1) & \cdots & x_n^{n-2}(x_n - x_1)
    \end{pmatrix}.
    \]
    Factor out $(x_i - x_1)$ from row $i$ for each $i \in \{2, \dots, n\}$:
    \[
    V_n(x_1, \dots, x_n) = \left(\prod_{i=2}^n (x_i - x_1)\right) \cdot V_{n-1}(x_2, \dots, x_n).
    \]
    By the induction hypothesis, $V_{n-1}(x_2, \dots, x_n) = \prod_{2 \le i < j \le n} (x_j - x_i)$.
    Combining these gives:
    \[
    V_n(x_1, \dots, x_n) = \left(\prod_{j=2}^n (x_j - x_1)\right) \left(\prod_{2 \le i < j \le n} (x_j - x_i)\right) = \prod_{1 \le i < j \le n} (x_j - x_i). \qedhere
    \]
\end{itemize}
\exinfo{Official Solution (Problem Sheet 17, Exercise 5). Proves the Vandermonde determinant formula by induction using column operations.}
\end{exercisesolution}

% Official Solution (Problem Sheet 17, Exercise 1)
\begin{exercisesolution}[Solution to \cref{exc:determinant_5x5_laplace}]
\label{sol:determinant_5x5_laplace}
Let $B = \begin{pmatrix}
  1 & 1 & 1 & 1 & 1 \\
  0 & 0 & 0 & 2 & 3 \\
  0 & 0 & 2 & 3 & 0 \\
  0 & 2 & 3 & 0 & 0 \\
  2 & 3 & 0 & 0 & 0
\end{pmatrix}$.
We expand $\det(B)$ along the first column:
\[
\det(B) = 1 \cdot (-1)^{1+1} \det \begin{pmatrix}
  0 & 0 & 2 & 3 \\
  0 & 2 & 3 & 0 \\
  2 & 3 & 0 & 0 \\
  3 & 0 & 0 & 0
\end{pmatrix} + 2 \cdot (-1)^{5+1} \det \begin{pmatrix}
  1 & 1 & 1 & 1 \\
  0 & 0 & 2 & 3 \\
  0 & 2 & 3 & 0 \\
  2 & 3 & 0 & 0
\end{pmatrix}.
\]
\begin{itemize}
    \item For the first $4 \times 4$ minor (anti-triangular), expanding along the bottom row gives:
    \[
    3 \cdot (-1)^{4+1} \det \begin{pmatrix} 0 & 2 & 3 \\ 2 & 3 & 0 \\ 3 & 0 & 0 \end{pmatrix} = -3 \cdot \left( 3 \cdot (-1)^{3+1} \det \begin{pmatrix} 2 & 3 \\ 3 & 0 \end{pmatrix} \right) = -9 \cdot (0 - 9) = 81.
    \]
    \item For the second $4 \times 4$ minor, performing $R_4 \to R_4 - 2R_1$ gives:
    \[
    \det \begin{pmatrix}
      1 & 1 & 1 & 1 \\
      0 & 0 & 2 & 3 \\
      0 & 2 & 3 & 0 \\
      0 & 1 & -2 & -2
    \end{pmatrix}
    = 1 \cdot \det \begin{pmatrix}
      0 & 2 & 3 \\
      2 & 3 & 0 \\
      1 & -2 & -2
    \end{pmatrix}
    = -2(-4 - 0) + 3(-4 - 3) = 8 - 21 = -13.
    \]
\end{itemize}
Thus:
\[
\det(B) = 1 \cdot 81 + 2 \cdot (-13) = 81 - 26 = 55.
\]
\begin{itemize}
    \item Over $\mathbb{R}$: $\det(B) = 55 \ne 0$, so $B$ is \textbf{invertible} over $\mathbb{R}$.
    \item Over $\mathbb{F}_5$: $\det(B) = 55 \equiv 0 \pmod 5$, so $B$ is \textbf{not invertible} over $\mathbb{F}_5$. \qedhere
\end{itemize}
\exinfo{Official Solution (Problem Sheet 17, Exercise 1). Computes $5 \times 5$ determinant via Laplace expansion and analyzes invertibility over $\mathbb{R}$ and $\mathbb{F}_5$.}
\end{exercisesolution}

% Official Solution (Problem Sheet 17, Exercise 6)
\begin{exercisesolution}[Solution to \cref{exc:tridiagonal_toeplitz_determinant}]
\label{sol:tridiagonal_toeplitz_determinant}
Let $D_n := \det(A_n)$. We verify the base cases:
\begin{itemize}
    \item For $n = 1$: $D_1 = \det(2) = 2 = 1 + 1$.
    \item For $n = 2$: $D_2 = \det \begin{pmatrix} 2 & 1 \\ 1 & 2 \end{pmatrix} = 4 - 1 = 3 = 2 + 1$.
\end{itemize}
For $n \ge 3$, expand $\det(A_n)$ along the first row:
\[
D_n = 2 \det(A_{n-1}) - 1 \cdot \det \begin{pmatrix}
  1 & 1 & 0 & \cdots & 0 \\
  0 & 2 & 1 & \cdots & 0 \\
  \vdots & 1 & 2 & \ddots & \vdots \\
  0 & 0 & \cdots & 1 & 2
\end{pmatrix}_{(n-1) \times (n-1)}.
\]
Expanding this second minor along its first column yields $1 \cdot \det(A_{n-2}) = D_{n-2}$.
Thus we obtain the linear recurrence relation:
\[
D_n = 2 D_{n-1} - D_{n-2} \iff D_n - D_{n-1} = D_{n-1} - D_{n-2}.
\]
The consecutive differences are constant: $D_n - D_{n-1} = D_2 - D_1 = 3 - 2 = 1$.
Therefore, $(D_n)_{n \ge 1}$ is an arithmetic progression:
\[
D_n = D_1 + (n-1) \cdot 1 = 2 + n - 1 = n + 1 \qquad \text{for all } n \ge 1. \qedhere
\]
\exinfo{Official Solution (Problem Sheet 17, Exercise 6). Derives and solves a second-order linear recurrence relation for tridiagonal Toeplitz determinants.}
\end{exercisesolution}

% Official Solution (Problem Sheet 17, Single Choice)
\begin{exercisesolution}[Solution to \cref{exc:parameter_for_3x3_determinant_sc}]
\label{sol:parameter_for_3x3_determinant_sc}
Let $A = \begin{pmatrix} 1 & 1 & 1 \\ 1 & x & 1 \\ 1 & 1 & x \end{pmatrix}$.
Applying row operations $R_2 \to R_2 - R_1$ and $R_3 \to R_3 - R_1$ gives the upper triangular form:
\[
\det(A) = \det \begin{pmatrix} 1 & 1 & 1 \\ 0 & x-1 & 0 \\ 0 & 0 & x-1 \end{pmatrix} = 1 \cdot (x-1) \cdot (x-1) = (x-1)^2.
\]
We solve the equation $(x-1)^2 = 1$:
\[
x - 1 = \pm 1 \implies x = 2 \quad \text{or} \quad x = 0.
\]
Among the given options, $x = 2$ is \textbf{(b)}.
\exinfo{Official Solution (Problem Sheet 17, Single Choice). Solves parametric $3 \times 3$ determinant equation via triangular row reduction.}
\end{exercisesolution}

% Official Solution (Problem Sheet 17, Exercise 7)
\begin{exercisesolution}[Solution to \cref{exc:properties_2x2_adjugate}]
\label{sol:properties_2x2_adjugate}
Let $A = \begin{pmatrix} A_{11} & A_{12} \\ A_{21} & A_{22} \end{pmatrix}$, so $\adj A = \begin{pmatrix} A_{22} & -A_{12} \\ -A_{21} & A_{11} \end{pmatrix}$ and $\det A = A_{11} A_{22} - A_{12} A_{21}$.

\begin{enumerate}[label=\textbf{(\alph*)}]
    \item \textbf{Adjugate identity:}
    \begin{align*}
    (\adj A) A &= \begin{pmatrix} A_{22} & -A_{12} \\ -A_{21} & A_{11} \end{pmatrix} \begin{pmatrix} A_{11} & A_{12} \\ A_{21} & A_{22} \end{pmatrix} \\
    &= \begin{pmatrix} A_{11}A_{22} - A_{12}A_{21} & A_{12}A_{22} - A_{12}A_{22} \\ -A_{11}A_{21} + A_{11}A_{21} & -A_{12}A_{21} + A_{11}A_{22} \end{pmatrix} \\
    &= \begin{pmatrix} \det A & 0 \\ 0 & \det A \end{pmatrix} = (\det A) I_2.
    \end{align*}
    Similarly, $A (\adj A) = (\det A) I_2$.

    \item \textbf{Determinant of the adjugate:}
    \[
    \det(\adj A) = A_{22} A_{11} - (-A_{12})(-A_{21}) = A_{11} A_{22} - A_{12} A_{21} = \det(A).
    \]

    \item \textbf{Transpose compatibility:}
    We have $A\transp = \begin{pmatrix} A_{11} & A_{21} \\ A_{12} & A_{22} \end{pmatrix}$.
    Applying the adjugate definition to $A\transp$:
    \[
    \adj(A\transp) = \begin{pmatrix} A_{22} & -A_{21} \\ -A_{12} & A_{11} \end{pmatrix} = \begin{pmatrix} A_{22} & -A_{12} \\ -A_{21} & A_{11} \end{pmatrix}\transp = (\adj A)\transp. \qedhere
    \]
\end{enumerate}
\exinfo{Official Solution (Problem Sheet 17, Exercise 7). Verifies the adjugate inverse identity, determinant formula, and transpose compatibility for $2 \times 2$ matrices.}
\end{exercisesolution}


\begin{ainote}
Determinants are built here in an order that looks circular and is not, and it
is worth saying once why. Existence comes from a recursion, the expansion along
a column, in which the sign of a permutation never appears; the Leibniz formula
and the sign are derived only afterwards, from the object already in hand.
\Cref{rem:no_circularity_sign} makes the point where it first arises.

Two hypotheses repay reading slowly. \Cref{prop:alternating_swap} assumes that
$D$ is $n$-linear and vanishes on matrices with two equal \emph{adjacent} rows,
and that adjacency is the whole strength of the statement: the two conditions
listed after it, vanishing whenever any two rows are equal and changing sign
under a swap, are conclusions drawn from that hypothesis rather than
alternative forms of it. And statement \textbf{(d)} of the properties of
$\det \colon \End(V) \to K$ on \cpageref{prop:det_endomorphism_properties}
needs $\dim V \ge 1$.

Prof. Biran marks three things \qt{exc.}: the column dichotomy inside
\cref{prop:inv_2x2}, the count $|S_n| = n!$, and the lower-triangular block
shape. All three are exercises here, solved above. The first is worth keeping in
his form rather than a weaker one. He writes that either column is a scalar
multiple of the other, which already covers the case of a zero column, and it is
that version the proof below him uses.
\end{ainote}

